\documentclass[11pt,a4paper]{article}

\usepackage[utf8]{inputenc}
\usepackage[T1]{fontenc}
\usepackage[ngerman]{babel}
\usepackage{amsmath,amssymb,amsthm}
\usepackage[a4paper,margin=2.5cm]{geometry}
\usepackage[hidelinks]{hyperref}
\usepackage{url}

\title{EABC-Operatoralgebra und Renormierungskette\\[0.3em]
	\large Ikosaeder--Dodekaeder-Dualität}
\author{Thomas Hoffbauer}
\date{}

\newtheorem{theorem}{Satz}
\newtheorem{definition}{Definition}
\newtheorem{remark}{Bemerkung}
\newtheorem{proposition}{Proposition}
\newtheorem{lemma}{Lemma}
\newtheorem{corollary}{Korollar}
\newtheorem{hypothesis}{Hypothese}
\newtheorem{conjecture}{Vermutung}

\newcommand{\R}{\mathbb{R}}
\newcommand{\Q}{\mathbb{Q}}
\newcommand{\Z}{\mathbb{Z}}

\begin{document}
	\maketitle
	
	\noindent\small
	\textbf{Reviewer Guide (1 page):}
	\texttt{REVIEWER\_GUIDE.pdf} --- what is proven, what is not, minimal hypotheses.
	\hrule
	\vspace{0.5em}
	
	\begin{abstract}
		\noindent
		(1)~Am Referenz-Ikosaeder gilt der isotrope geometrische Fixpunkt
		$M_{\mathrm{geom}}(I_{12},\sigma_{\mathrm{geo}})=24I_3$.
		(2)~Eine EABC-Primzahl $p$ erzeugt den effektiven Rang-1-Defekt
		$M_{\mathrm{eff}}(K^+)=24I_3+w_pvv^{\mathsf T}$ mit $\|v\|=1$ und messbarer
		Anisotropie $\Delta(M_{\mathrm{eff}}(K^+))=w_p$.
		(3)~Der Retraktionsoperator $R^*_{\mathrm{EABC}}$ entfernt diesen Defekt projektiv.
		(4)~Es folgt die \emph{24-Isotropy-Restoration}:
		$M_{\mathrm{eff}}(R^*_{\mathrm{EABC}}(K^+))=24I_3$.
		(5)~Die Satzkette ist auf dem Peano-Schalenstapel global formalisiert
		(Lean~4.29, \texttt{sorry}\,$=\,0$).
	\end{abstract}
	
	%=============================================================================
	\section{Einleitung: historische und moderne Einordnung}
	\label{sec:introduction}
	%=============================================================================
	
	Die EABC-Renormierungskette verbindet diskrete Arithmetik, polyedrische Symmetrie
	und effektive Tensorphysik in einem einzigen, formal prüfbaren Modell.
	Mehrere historische Linien konvergieren hier --- ohne die Resultate dieser Arbeit
	als deren unmittelbare Fortsetzung zu beanspruchen.
	
	\paragraph{Leonhard Euler (1707--1783).}
	Eulers Polyederformel und seine systematische Theorie der platonischen Körper
	verankern das Ikosaeder als \emph{mathematisches} Referenzobjekt, nicht bloß als
	Zeichnung.
	Die zwölf Ecken des regulären Ikosaeders auf $S^2$ und die zugehörige
	$A_5$-Symmetrie sind der geometrische Boden von $M_{\mathrm{geom}}$;
	Theorem~\ref{thm:geom-fixpoint} ist in diesem Sinne ein Euler-artiger Satz:
	eine endliche, explizite Summe symmetrischer dyadischer Produkte kollabiert
	zu einem skalaren Vielfachen von $I_3$.
	Die Peano-Schalenstapel (Abschnitt~\ref{sec:formal-verification}) erinnern an
	Eulers Methode, diskrete Strukturen schichtweise über $\mathbb{N}$ zu erzeugen.
	
	\paragraph{Bernhard Riemann (1826--1866).}
	Riemanns Trennung von lokalen Koordinaten und globaler Struktur auf Mannigfaltigkeiten
	entspricht methodisch der hier durchgeführten Separation von
	$M_{\mathrm{eff}}$ (Konfigurationsraum) und $M_{\mathrm{geom}}$ (geometrischer Fixpunkt,
	Abschnitt~\ref{sec:two-levels}).
	Der Ordnungsparameter $\Delta(M)=\lambda_{\max}-\lambda_{\min}$ ist ein
	spektral-geometrisches Funktional im riemannschen Sinne: Isotropie bedeutet
	Entartung des Spektrums.
	Die Primzahlen $p=2n+1$ treten als \emph{diskrete Normschalen-Indizes} auf ---
	nicht als analytische Objekte der Riemannschen $\zeta$-Funktion.
	Der strukturelle Bezug zu Riemann liegt in der Frage, wie globale Invarianz
	($24I_3$) aus lokaler arithmetischer Störung ($w_p$) durch Projektion
	wiederhergestellt wird.
	
	\paragraph{Paul Dirac (1902--1984).}
	Diracs Vakuum-Bild --- ein symmetrischer Grundzustand, dem diskrete
	Anregungen (Massen) als Störungen überlagert sind --- spiegelt sich in der
	Zwei-Ebenen-Architektur wider:
	$M_{\mathrm{geom}}=24I_3$ ist das isotrope ``Vakuum'',
	$M_{\mathrm{eff}}(K^+)=24I_3+w_pvv^{\mathsf T}$ die messbare Anisotropie
	(Theorem~\ref{thm:prime-delta}, \texttt{rankOne\_defect\_not\_isotropic}).
	Die Rang-1-Störung $w_pvv^{\mathsf T}$ ist die minimalste nicht-triviale
	Symmetriebrechung in $\mathrm{Sym}_3(\mathbb{R})$; $\Delta(M^+)=w_p$ quantisiert
	das ``Massen-Residuum'' der Störung.
	Dirac würde die Retraktion $R^*$ nicht als Dynamik, sondern als
	Zustandsreduktion lesen: Rückkehr zum symmetrischen Referenzzustand.
	
	\paragraph{Gerard 't~Hooft (*1946).}
	Die Renormierungs- und Eichtheorie-Tradition, deren Höhepunkt 't~Hoofts Arbeit
	zur renormierbaren nicht-abelschen Eichfeldtheorie bildet, liefert das
	konzeptuelle Vokabular des \emph{Retraktionsoperators} $R^*_{\mathrm{EABC}}$:
	kein ad-hoc-Glätten, sondern projektive Rückführung auf einen wohldefinierten
	Fixpunkt-Teilraum $\mathcal{F}$.
	Die Idempotenz $(R^*)^2=R^*$ (Lemma~\ref{lem:idempotence}) entspricht dem
	Gedanken, dass ein korrekt gewählter ``Gauge-Fix'' nach einem Schritt stabil bleibt.
	Die EABC-Theorie beansprucht keine QFT-Renormierungsgruppe im analytischen Sinn;
	sie übernimmt die \emph{Struktur}: lokaler Defekt, globaler Fixpunkt, ein Schritt genügt.
	
	\paragraph{Zeitgenössische Nachbarn.}
	Maynard, Tao und Scholze behandeln Primzahlen auf anderem Niveau
	(Gaps, additive Kombinatorik, perfektoide Geometrie; Abschnitt~\ref{sec:modern-neighbors}).
	Die EABC-Theorie bleibt diskret-modular: $p\bmod 12$ klassifiziert Normschalen,
	nicht die Primzahlverteilung selbst.
	
	%=============================================================================
	\section{EABC-Familien und Anisotropietensor}
	%=============================================================================
	
	Die vier invertierbaren Restklassen modulo~$12$ definieren die EABC-Familien
	\[
	E \equiv 1,\qquad A \equiv 5,\qquad B \equiv 7,\qquad C \equiv 11 \pmod{12},
	\]
	mit Gewichten $w_E=1$, $w_A=5$, $w_B=7$, $w_C=11$.
	
	Sei $\{v_i\}_{i=1}^{12}\subset S^2\subset\R^3$ ein 12-Kuss-System (Einheitsrichtungen)
	und $\sigma\colon\{1,\ldots,12\}\to\{E,A,B,C\}$ eine Label-Zuweisung.
	Der \emph{Anisotropietensor} ist
	\[
	M(\sigma)
	\;=\;
	\sum_{i=1}^{12} w_{\sigma(i)}\, v_i v_i^{\mathsf T}
	\;\in\;
	\mathrm{Sym}_3(\R).
	\]
	
	\begin{definition}[Ordnungsparameter]
		\[
		\Delta(M) \;:=\; \lambda_{\max}(M) - \lambda_{\min}(M),
		\]
		wobei $\lambda_1\le\lambda_2\le\lambda_3$ die Eigenwerte von $M$ bezeichnen.
		Isotropie entspricht $\Delta(M)=0$.
	\end{definition}
	
	%=============================================================================
	\section{Renormierungsoperatoren}
	%=============================================================================
	
	Sei $K_n=(V_n,\sigma_n)$ eine Konfiguration aus Vertices $V_n$ und Zuweisung $\sigma_n$.
	
	\begin{definition}[$R_{\mathrm{EABC}}$]
		Der Operator $R_{\mathrm{EABC}}$ projiziert $\sigma$ auf die geometrisch zulässige
		3+3+3+3-Zuweisung: pro Ecke $v_i$ wird die mod-$12$-Restklasse
		$r(v_i)$ aus den skalierten Koordinaten und die nächste EABC-Restklasse
		bestimmt. $R_{\mathrm{EABC}}$ ist idempotent; Fixpunkte erfordern perfekte
		geometrische Konsistenz der Koordinaten-Klassifikation.
	\end{definition}
	
	\begin{definition}[$P_I$ (Voronoi-Pullback)]
		$P_I$ bildet gestörte Kussrichtungen auf das Referenz-Ikosaeder ab:
		\[
		P_I(v) \;=\; \arg\max_{r_j\in I_{12}} \langle v, r_j\rangle,
		\]
		wobei $I_{12}$ die 12 Referenz-Ecken des regulären Ikosaeders auf $S^2$ ist.
	\end{definition}
	
	\begin{definition}[$R^*_{\mathrm{EABC}}$]
		Der robuste Fixpunktoperator ist die Komposition
		\[
		R^*_{\mathrm{EABC}} \;:=\; P_I \circ R_{\mathrm{EABC}}.
		\]
		Zuerst geometrische Rückführung, dann Label-Projektion.
	\end{definition}
	
	\subsection{Konfigurationsraum und Retraktion}
	
	\begin{definition}[Konfigurationsraum $\mathcal{C}$]
		$\mathcal{C}$ ist die Menge aller Konfigurationen
		\[
		K=(V,\sigma),
		\qquad
		V=(v_1,\ldots,v_{12})\subset S^2,
		\qquad
		\sigma\colon\{1,\ldots,12\}\to\{E,A,B,C\},
		\]
		also aller (ggf.\ gestörten) EABC-Kuss-Konfigurationen.
	\end{definition}
	
	\begin{definition}[Fixpunkt-Teilraum $\mathcal{F}$]
		\[
		\mathcal{F}
		\;:=\;
		\bigl\{K\in\mathcal{C} : R^*_{\mathrm{EABC}}(K)=K\bigr\}
		\;\subset\;
		\mathcal{C}
		\]
		ist der Teilraum der \emph{robusten Fixpunkte} des Operators $R^*_{\mathrm{EABC}}$.
		Das geometrische Referenzobjekt $I_{12}^{\mathrm{EABC}}$ liegt in $\mathcal{F}$.
	\end{definition}
	
	\begin{definition}[Einbettung und Retraktion]
		Die kanonische Einbettung ist
		\[
		\iota\colon \mathcal{F}\hookrightarrow \mathcal{C},\qquad \iota(K)=K.
		\]
		Der robuste Operator wirkt als Abbildung
		\[
		R^*_{\mathrm{EABC}}\colon \mathcal{C}\longrightarrow \mathcal{F}.
		\]
	\end{definition}
	
	\begin{proposition}[Retraktion auf $\mathcal{F}$]
		\label{prop:retraction}
		Mit $\iota$ und $R^*_{\mathrm{EABC}}$ wie oben gilt:
		\begin{enumerate}
			\item[(a)] \textbf{Idempotenz auf $\mathcal{C}$:}
			$(R^*_{\mathrm{EABC}})^2 = R^*_{\mathrm{EABC}}$.
			\item[(b)] \textbf{Fixpunkt auf $\mathcal{F}$:}
			$R^*_{\mathrm{EABC}}(K)=K$ für alle $K\in\mathcal{F}$.
			\item[(c)] \textbf{Retraktionsidentität:}
			$R^*_{\mathrm{EABC}}\circ \iota = \mathrm{id}_{\mathcal{F}}$.
		\end{enumerate}
		Insbesondere ist $R^*_{\mathrm{EABC}}\colon\mathcal{C}\to\mathcal{F}$ eine
		\emph{Retraktion} von $\mathcal{C}$ auf den Fixpunkt-Teilraum $\mathcal{F}$.
	\end{proposition}
	
	\begin{proof}[Skizze]
		(a) ist Lemma~\ref{lem:idempotence} unten.
		Für (b) folgt aus der Definition von $\mathcal{F}$.
		Für (c) setze $K\in\mathcal{F}$; dann ist
		$R^*_{\mathrm{EABC}}(\iota(K))=R^*_{\mathrm{EABC}}(K)=K=\iota(K)$.
	\end{proof}
	
	\begin{proposition}[Konvergenz in einem Schritt]
		\label{prop:convergence}
		Ist $R^*_{\mathrm{EABC}}$ idempotent, so gilt für alle $K\in\mathcal{C}$ und
		alle $n\ge 1$:
		\[
		(R^*_{\mathrm{EABC}})^n(K) \;=\; R^*_{\mathrm{EABC}}(K).
		\]
		Insbesondere konvergiert die Renormierungsiteration nach \emph{einem} Schritt
		auf das Bild $\mathcal{F}$.
	\end{proposition}
	
	\begin{proof}
		Vollständig formalisiert in \texttt{eabc-renorm/EabcRenorm/Retraction.lean}
		als \texttt{IsIdempotent.iterate\_ge\_one}.
		Für $n=1$ trivial; für $n\ge 2$ folgt $(R^*)^n=(R^*)^{n-1}\circ R^*$
		und Induktion über die Idempotenz $(R^*)^2=R^*$.
	\end{proof}
	
	%=============================================================================
	\section{Fixpunkt und Defektklassen}
	%=============================================================================
	
	\begin{definition}[Geometrischer Fixpunkt $I_{12}^{\mathrm{EABC}}$]
		Die Konfiguration $I_{12}^{\mathrm{EABC}}$ besteht aus den 12 Ikosaeder-Ecken
		mit der geometrischen EABC-Zuweisung (je 3 Labels pro Familie).
	\end{definition}
	
	\begin{lemma}[Lemma 1A — Spur am balancierten Fixpunkt]
		\label{lem:trace}
		Sei $M(\sigma)=\sum_{i=1}^{12} w_{\sigma(i)}\, v_i v_i^{\mathsf T}$ mit $\|v_i\|=1$
		und 3+3+3+3-Balance der Label-Zuweisung $\sigma$.
		Dann gilt $\operatorname{tr}(M)=72$.
	\end{lemma}
	
	\begin{proof}
		Da $\operatorname{tr}(v_i v_i^{\mathsf T})=\|v_i\|^2=1$, folgt
		\[
		\operatorname{tr}(M)
		\;=\;
		\sum_{i=1}^{12} w_{\sigma(i)}
		\;=\;
		3(1+5+7+11)
		\;=\;
		72.
		\]
	\end{proof}
	
	\begin{lemma}[Lemma 1B — Skalar aus Isotropie]
		\label{lem:scalar}
		Ist $M\in\mathrm{Sym}_3(\R)$ isotrop, d.\,h.\ $M=\lambda I$ für ein $\lambda\in\R$,
		und gilt $\operatorname{tr}(M)=72$, dann folgt $\lambda=24$ und
		$\lambda_1=\lambda_2=\lambda_3=24$.
	\end{lemma}
	
	\begin{proof}
		Aus $M=\lambda I$ folgt $\operatorname{tr}(M)=3\lambda=72$, also $\lambda=24$.
	\end{proof}
	
	\begin{lemma}[Lemma 1C — Geometrische Zuordnung der mod-12-Bahnen]
		\label{lem:orbit-assignment}
		Die Ikosaeder-Ecken zerfallen bei Skala $s=6$ in vier Bahnen
		$\mathcal{O}_2,\mathcal{O}_4,\mathcal{O}_8,\mathcal{O}_{10}$
		(je drei Ecken).
		Die Abbildung \emph{nächste EABC-Restklasse} liefert
		\[
		\mathcal{O}_2\mapsto E\;(w=1),\quad
		\mathcal{O}_4\mapsto A\;(w=5),\quad
		\mathcal{O}_8\mapsto B\;(w=7),\quad
		\mathcal{O}_{10}\mapsto C\;(w=11).
		\]
		Diese Zuordnung ist injektiv auf den Familien und fest durch
		$r(v)=(\mathrm{round}(sx)+\mathrm{round}(sy)+\mathrm{round}(sz))\bmod 12$
		und die Kreis-Distanz auf $\Z/12\Z$.
	\end{lemma}
	
	\begin{proof}
		Die vier Aussagen $\mathcal{O}_r\mapsto\{\text{E,A,B,C}\}$ sind endliche
		Verifikationen auf den Restklassen $2,4,8,10$; formalisiert in
		\texttt{eabc-renorm/EabcRenorm/Mod12.lean} per \texttt{native\_decide}.
	\end{proof}
	
	\begin{lemma}[Lemma 1D — Geometrische Isotropie]
		\label{lem:geom-iso}
		Am geometrischen Fixpunkt $I_{12}^{\mathrm{EABC}}$ gilt
		\[
		M(\sigma)
		\;=\;
		\sum_{i=1}^{12} w_{\sigma(i)}\, v_i v_i^{\mathsf T}
		\;=\;
		24\,I_3.
		\]
	\end{lemma}
	
	\begin{proof}[Zwei unabhängige Wege]
		\textbf{Weg I (direkte Summe).}
		Mit den expliziten Koordinaten aus Anhang~\ref{app:generators} ergibt die
		Summe der zwölf gewichteten äußeren Produkte in $\Q(\sqrt5)$ exakt $24I_3$.
		
		\smallskip\noindent
		\textbf{Weg II (Symmetrie + Schur).}
		Sei $\mathcal{I}\cong A_5$ die Ikosaeder-Rotationsgruppe.
		Die gewichtete Bahn-Summe $M=\sum_r w_{\phi(r)}T_r$ ist $\mathcal{I}$-invariant;
		nach Schur folgt $M=\lambda I$.
		Lemma~\ref{lem:trace} liefert $\lambda=24$.
		Die Erzeuger-Verifikation steht in Anhang~\ref{app:generators}.
	\end{proof}
	
	\begin{theorem}[Isotropie-Lemma — Komposition]
		\label{thm:isotropy}
		Am geometrischen EABC-Fixpunkt $I_{12}^{\mathrm{EABC}}$ hat $M(\sigma)$ die
		Eigenwerte $\lambda_1=\lambda_2=\lambda_3=24$, also $\Delta(M)=0$.
	\end{theorem}
	
	\begin{proof}
		Lemma~\ref{lem:geom-iso} liefert $M=24I_3$ direkt.
		Alternativ: Lemma~\ref{lem:geom-iso} (Weg II) gibt Isotropie,
		Lemma~\ref{lem:trace} gibt $\operatorname{tr}(M)=72$,
		Lemma~\ref{lem:scalar} gibt $\lambda=24$.
	\end{proof}
	
	\begin{remark}[Lean-Formalisierung (V4)]
		In \texttt{eabc-renorm} ist der Beweis dreistufig angelegt:
		\texttt{trace\_geometric\_fixpoint} (1A),
		\texttt{isotropy\_from\_trace} (1B),
		\texttt{geometric\_fixpoint\_isotropy} (1C/1D).
		Stand: \texttt{lake build} grün, \texttt{sorry}\,$=\,0$
		(siehe Abschnitt~\ref{sec:formal-verification}).
	\end{remark}
	
	\begin{remark}[Numerische Bestätigung]
		Unabhängig bestätigt \texttt{eabc\_icosahedron\_test.py},
		\texttt{run\_core\_lemma\_tests()} die Aussage (4/4 PASS, Toleranz $10^{-6}$).
	\end{remark}
	
	\begin{lemma}[Projektor-Idempotenz]
		\label{lem:idempotence}
		$R^*_{\mathrm{EABC}} = P_I \circ R_{\mathrm{EABC}}$ ist idempotent:
		\[
		(R^*_{\mathrm{EABC}})^2 \;=\; R^*_{\mathrm{EABC}}.
		\]
	\end{lemma}
	
	\begin{proof}[Skizze]
		Nach $P_I$ liegen alle Vertices auf den Referenz-Ecken; $R_{\mathrm{EABC}}$ ist
		auf dieser Untervarietät bereits Fixpunktabbildung. Ein zweiter Schritt ändert
		weder Geometrie noch Labels.
	\end{proof}
	
	\smallskip\noindent
	\textit{Numerisch verifiziert in \texttt{eabc\_icosahedron\_test.py},
		\texttt{run\_core\_lemma\_tests()} (4/4 PASS, Toleranz $10^{-6}$).}
	
	\begin{definition}[Defektklassen]
		\begin{align*}
			\text{FIXPOINT}            &\colon \Delta(M)\approx 0,\; 3{+}3{+}3{+}3\text{-Balance},\; \text{ungestörte Geometrie},\\
			\text{LABEL\_DEFECT}       &\colon \text{Label-Zuweisung ohne 3+3+3+3-Balance (vor } R_{\mathrm{EABC}}\text{)},\\
			\text{DIRECTION\_DEFECT}   &\colon \text{balancierte Labels, aber }\Delta(M)>\varepsilon \text{ bei gestörter Geometrie},\\
			\text{CLASSIFICATION\_DEFECT} &\colon \text{mod-}12\text{-Balance der geometrischen Zuweisung gebrochen},\\
			\text{PRIME\_DEFECT}(p)    &\colon p=2n+1\in\mathbb{P},\; p\bmod 12\in\{1,5,7,11\},\;
			\text{neue Prim-Normkugel verletzt zunächst die 3+3+3+3-Balance}.
		\end{align*}
	\end{definition}
	
	%=============================================================================
	\section{Prim-Erzeugungsschritt}
	\label{sec:prime-step}
	%=============================================================================
	
	\subsection{Operation $2n\mapsto 2n+1$}
	
	Sei $N=2n$ eine gerade Normschale.
	Der \emph{Prim-Erzeugungsschritt} bildet
	\[
	P(2n)
	\;:=\;
	\begin{cases}
		2n+1, & \text{falls }2n+1\in\mathbb{P},\\
		\bot, & \text{sonst}.
	\end{cases}
	\]
	Für $p=P(2n)$ gilt bei Primzahlen $p>3$ stets
	$p\bmod 12\in\{1,5,7,11\}$; die neue Normkugel landet damit in genau
	einer EABC-Familie.
	
	\begin{definition}[Erweiterung um eine Prim-Normkugel]
		Gegeben $K_n\in\mathcal{F}$ mit Label-Zuweisung $\sigma$ und $p=P(2n)$.
		Die erweiterte Konfiguration $K_n^+$ entsteht durch Hinzufügen einer
		Einheitsrichtung $v\in S^2$ mit Label $\mathrm{EABC}(p)$:
		\[
		M^+ \;=\; M + w_p\, v v^{\mathsf T},
		\qquad
		w_p\in\{1,5,7,11\}\text{ gemäß }p\bmod 12.
		\]
		Die Familienbesetzung wird von $(3,3,3,3)$ zu $(4,3,3,3)$ (in genau einer Komponente).
	\end{definition}
	
	\begin{proposition}[PRIME\_DEFECT $\Rightarrow$ LABEL\_DEFECT]
		\label{prop:prime-implies-label}
		Ist $\text{PRIME\_DEFECT}(p)$ aktiv, so gilt $K_n^+\notin\mathcal{F}$ und
		$R^*_{\mathrm{EABC}}(K_n^+)\neq K_n^+$.
		Die Störung ist ein \emph{arithmetisch privilegierter} Label-Defekt.
	\end{proposition}
	
	\subsection{Anisotropie nach dem Primschritt}
	
	Am Fixpunkt war $M=24I_3$.
	Nach dem Primschritt gilt im Rang-1-Modell
	\[
	M^+ \;=\; 24\,I_3 + w_p\, v v^{\mathsf T},
	\qquad \|v\|=1.
	\]
	
	\begin{theorem}[Primdefekt-Satz]
		\label{thm:prime-delta}
		Für $w_p>0$ und $\|v\|=1$ hat $M^+=24I_3+w_p vv^{\mathsf T}$ die Eigenwerte
		$\lambda_1=\lambda_2=24$ und $\lambda_3=24+w_p$ (bzw.\ Permutation falls $v$ nicht
		Koordinatenachse ist).
		Insbesondere
		\[
		\Delta(M^+) \;=\; w_p \;\in\; \{1,5,7,11\}.
		\]
	\end{theorem}
	
	\begin{proof}
		Für symmetrische Rang-1-Störungen $w\,vv^{\mathsf T}$ mit $\|v\|=1$ ist $v$ Eigenvektor
		zum Eigenwert $w$; die Orthogonalen zu $v$ haben Eigenwert $0$.
		Also hat $24I_3+w_p vv^{\mathsf T}$ die Eigenwerte $24$ (doppelt) und $24+w_p$ (einfach).
		Damit $\Delta(M^+)=\lambda_{\max}-\lambda_{\min}=w_p$.
		Formalisiert in \texttt{eabc-renorm/EabcRenorm/PrimeDefect.lean}
		als \texttt{delta\_rank\_one\_perturbation}.
	\end{proof}
	
	\begin{remark}[Primzahlen als Symmetriebrecher]
		Eine Primzahl erscheint hier nicht als bloße Zahl ohne Teiler, sondern als
		\emph{neue ungerade Normkugel}, die aus einer geraden Schale hervorgeht,
		ein zulässiges EABC-Label trägt, die isotrope Balance bricht und einen
		Renormierungsimpuls auslöst.
	\end{remark}
	
	\subsection{$R^*_{\mathrm{EABC}}$ als Prim-Entstörungsoperator}
	
	Bei einem Primschritt hat $R^*_{\mathrm{EABC}}=P_I\circ R_{\mathrm{EABC}}$ die Bedeutung
	\[
	\text{Prim-Erzeugung}
	\;\Rightarrow\;
	\text{PRIME\_DEFECT}
	\;\Rightarrow\;
	\text{projektive Rückführung}.
	\]
	Die Kette lautet
	\[
	K_n \;\longrightarrow\; K_n^+ \;\longrightarrow\; R^*_{\mathrm{EABC}}(K_n^+).
	\]
	Nach Proposition~\ref{prop:convergence} genügt \emph{ein} Projektionsschritt:
	$(R^*_{\mathrm{EABC}})^2=R^*_{\mathrm{EABC}}$.
	
	\begin{theorem}[Prim-Normkugel-Satz — Label-Teil]
		\label{thm:prime-norm}
		Sei $K_n\in\mathcal{F}$ mit geometrischer Label-Zuweisung und
		$K_n^+=\mathrm{Extend}(p)\,K_n$ für $p=P(2n)\in\mathbb{P}$, $p>3$.
		Dann gilt:
		\begin{enumerate}
			\item[(a)] $\text{PRIME\_DEFECT}(p)$ und $K_n^+\notin\mathcal{F}$;
			\item[(b)] $R^*_{\mathrm{EABC}}(K_n^+)\in\mathcal{F}$ mit wiederhergestellter
			3+3+3+3-Balance;
			\item[(c)] ein Schritt genügt: $(R^*_{\mathrm{EABC}})^n(K_n^+)=R^*_{\mathrm{EABC}}(K_n^+)$
			für alle $n\ge 1$.
		\end{enumerate}
	\end{theorem}
	
	\begin{proof}[Beweisidee]
		\textbf{(a)} Die virtuelle Prim-Normkugel erhöht die Multiketten-Zählung einer
		Familie von $3$ auf $4$ bei unveränderter 12er-Geometrie
		(Lemma~\ref{lem:prime-imbalance}, formalisiert als
		\texttt{extendPrime\_creates\_defect}).
		
		\smallskip\noindent
		\textbf{(b)} Entscheidend ist, dass $R_{\mathrm{EABC}}$ \emph{labelblind} ist:
		Es setzt die Zuweisung unabhängig von der Eingabe auf die geometrische
		3+3+3+3-Struktur und löscht den Prim-Schalen-Defekt
		(\texttt{rEabc\_restores\_balance} in \texttt{PrimeDefect.lean}).
		Nach $P_I$ liegt die Geometrie wieder auf dem Referenz-Ikosaeder.
		
		\smallskip\noindent
		\textbf{(c)} Folgt aus der Idempotenz von $R^*_{\mathrm{EABC}}$
		(Proposition~\ref{prop:convergence}).
	\end{proof}
	
	\subsection{Die zwei Tensorebenen: Geometrischer Fixpunkt versus effektiver Konfigurationsraum}
	\label{sec:two-levels}
	
	Ein zentrales Missverständnis bei diskreten Symmetriebrechungen liegt in der
	unzulässigen Vermischung von physikalischem Zustand und geometrischem Hintergrund.
	Das EABC-Renormierungsmodell löst diesen Konflikt durch eine strikte Separation
	zweier Tensorebenen --- auf jeder Peano-Schale $n\in\mathbb{N}$ und zunächst
	unabhängig von metrischen Radien (Phase~A in Abschnitt~\ref{sec:formal-verification}).
	
	Der Primdefekt wirkt gleichzeitig auf Label- und Tensorebene:
	\begin{enumerate}
		\item[(L)] Label-Multikette: $(4,3,3,3)$ statt $(3,3,3,3)$;
		\item[(T)] Tensor: $M_{\mathrm{eff}}(K^+)=24I_3+w_pvv^{\mathsf T}$ (Rang-1-Störung).
	\end{enumerate}
	
	\subsubsection{Die rein geometrische Ebene: der Fixpunkt $M_{\mathrm{geom}}$}
	
	Die Basis bildet der rein geometrische Tensor $M_{\mathrm{geom}}$, der ausschließlich
	die ungestörte Ikosaeder-Konfiguration abbildet.
	Für Vertices $V=(v_i)_{i=1}^{12}$ und Label-Zuweisung $\sigma$ gilt
	(Definition~\ref{def:m-geom}, \texttt{M\_geom} in \texttt{EffectiveTensor.lean})
	\[
	M_{\mathrm{geom}}(V,\sigma)
	\;=\;
	\sum_{i=1}^{12} w_{\sigma(i)}\, v_i v_i^{\mathsf T}.
	\]
	Das in \texttt{IcosahedronSymmetry.lean} verifizierte Theorem
	\texttt{geometric\_fixpoint\_isotropy} beweist, dass für das Referenzobjekt
	$I_{12}^{\mathrm{EABC}}$ mit $\sigma_{\mathrm{geo}}$ alle Richtungsabhängigkeiten
	exakt kollabieren:
	\[
	\boxed{
		M_{\mathrm{geom}}(I_{12},\sigma_{\mathrm{geo}}) \;=\; 24I_3
	}.
	\]
	Dieser Zustand ist perfekt isotrop und invariant unter $SO(3)$.
	Er enthält keine dynamischen Defektfelder; in Lean entspricht er
	\texttt{geometricM} bzw.\ \texttt{M\_geom\_geometric\_fixpoint}.
	
	\subsubsection{Die effektive physikalische Ebene: der Konfigurationsraum $M_{\mathrm{eff}}$}
	
	Der effektive Tensor $M_{\mathrm{eff}}$ operiert auf dem Konfigurationsraum
	$\mathcal{C}$ (Lean: \texttt{Config} mit Feldern \texttt{label},
	\texttt{primeShell}, \texttt{tensorDefect}).
	Tritt ein EABC-Primereignis $p$ mit Gewicht $w_p\in\{1,5,7,11\}$ entlang eines
	normierten Defektvektors $v$ ($\|v\|=1$) auf, so setzt man
	(Definition~\ref{def:m-eff})
	\[
	K^+ \;:=\; \mathrm{Extend}(p,v)\,K,
	\qquad
	M_{\mathrm{eff}}(K^+) \;=\; 24I_3 + w_p\,vv^{\mathsf T}.
	\]
	In Lean: \texttt{extendPrimeTensor} schreibt \texttt{tensorDefect = some (w\_p, v)},
	und \texttt{M\_eff\_with\_defect} liefert \texttt{rankOnePerturbation 24 w v}.
	Das Theorem \texttt{rankOne\_defect\_not\_isotropic} zeigt: für $w_p>0$ ist
	$M_{\mathrm{eff}}(K^+)$ nicht isotrop.
	Die messbare Anisotropie ist
	\[
	\Delta\bigl(M_{\mathrm{eff}}(K^+)\bigr) \;=\; w_p
	\]
	(Theorem~\ref{thm:prime-delta}, \texttt{delta\_prime\_defect\_at\_fixpoint}).
	
	\subsubsection{Die Wirkung von $R^*_{\mathrm{EABC}}$: ein reiner Projektionssatz}
	
	$R^*_{\mathrm{EABC}}$ ist kein Dämpfungs- oder Glättungsoperator, sondern ein
	algebraischer Retrakt auf $\mathcal{F}$.
	Die Tensor-Restauration geschieht nicht analytisch \emph{innerhalb} der gestörten
	Matrix $M_{\mathrm{eff}}(K^+)$, sondern durch projektiven Kollaps der Konfiguration
	zurück auf die geometrische Referenz:
	\[
	K^+
	\;\xrightarrow{\;R^*_{\mathrm{EABC}}\;}
	(I_{12},\sigma_{\mathrm{geo}})
	\;\xrightarrow{\;M_{\mathrm{geom}}\;}
	24I_3.
	\]
	Formal faktorisiert die Kette in zwei unabhängige Schritte:
	\begin{enumerate}
		\item \textbf{Projektion (Ebene~1):}
		$M_{\mathrm{eff}}(R^*(K))=M_{\mathrm{geom}}(I_{12},\sigma_{\mathrm{geo}})$
		ist \emph{definitorisch}
		(Lemma~\ref{lem:tensor-projection}, \texttt{M\_eff\_projection}:
		\texttt{rStar} setzt \texttt{tensorDefect} auf \texttt{none}).
		\item \textbf{Fixpunkt (Ebene~2):}
		$M_{\mathrm{geom}}(I_{12},\sigma_{\mathrm{geo}})=24I_3$
		ist der \emph{inhaltliche} Ikosaeder-Satz
		(Theorem~\ref{thm:geom-fixpoint}).
	\end{enumerate}
	Für eine einzelne Schale liefert \texttt{prime\_norm\_full\_restoration} beides:
	$R^*(K^+)\in\mathcal{F}$ und $M_{\mathrm{eff}}(R^*(K^+))=24I_3$.
	Auf einem Peano-Schalenstapel gilt dasselbe komponentenweise
	(\texttt{all\_shells\_tensor\_restored} in \texttt{ShellStack.lean}).
	
	\begin{table}[h!]
		\centering
		\small
		\begin{tabular}{p{0.14\linewidth}p{0.36\linewidth}p{0.22\linewidth}p{0.22\linewidth}}
			\hline
			\textbf{Ebene} & \textbf{Mathematisches Objekt} & \textbf{Symmetrie} & \textbf{Lean} \\
			\hline
			Geometrisch &
			$M_{\mathrm{geom}}(I_{12},\sigma_{\mathrm{geo}})=24I_3$ &
			Isotropie (Referenz) &
			\texttt{geometric\_fixpoint\_isotropy} \\
			Physikalisch &
			$M_{\mathrm{eff}}(K^+)=24I_3+w_pvv^{\mathsf T}$ &
			Rang-1-Anisotropie &
			\texttt{extendPrimeTensor}, \texttt{rankOne\_defect\_not\_isotropic} \\
			Projektiv &
			$M_{\mathrm{eff}}(R^*(K^+))=24I_3$ &
			restaurierte Isotropie &
			\texttt{prime\_norm\_full\_restoration} \\
			Peano (global) &
			$\mathcal{R}^*(\mathcal{S}_{\mathrm{def}})=\mathcal{S}_{\mathrm{iso}}$ &
			schalenweise unabhängig &
			\texttt{all\_shells\_tensor\_restored} \\
			\hline
		\end{tabular}
		\caption{Gegenüberstellung geometrischer Invarianz, effektiver Störung und projektiver Restauration.
			\texttt{IsEabcFixpoint} beschreibt den Fixpunkt $K\in\mathcal{F}$, nicht direkt den Tensor;
			die Isotropie $24I_3$ folgt über \texttt{M\_eff} und Theorem~\ref{thm:geom-fixpoint}.}
		\label{tab:two-levels}
	\end{table}
	
	\subsection{Tensor-Restauration als Projektionssatz}
	\label{sec:tensor-projection}
	
	Der gesamte inhaltliche Kern von Phase~A lässt sich in einem Flussdiagramm
	zusammenfassen (Hauptsatz: Theorem~\ref{thm:tensor-restore}):
	
	\begin{center}
		\fbox{%
			\begin{minipage}{0.55\linewidth}
				\centering
				\vspace{2mm}
				$M_{\mathrm{geom}}(I_{12},\sigma_{\mathrm{geo}}) = 24I_3$ \\[6pt]
				$\downarrow$ \\[6pt]
				$M_{\mathrm{eff}}(K^+) = 24I_3 + w_p\,vv^{\mathsf T}$ \\[6pt]
				$\downarrow$ \\[6pt]
				$\Delta\bigl(M_{\mathrm{eff}}(K^+)\bigr) = w_p$ \\[6pt]
				$\downarrow$ \\[6pt]
				$R^*_{\mathrm{EABC}}$ \\[6pt]
				$\downarrow$ \\[6pt]
				$\boxed{M_{\mathrm{eff}}\bigl(R^*_{\mathrm{EABC}}(K^+)\bigr) = 24I_3}$
				\vspace{2mm}
			\end{minipage}%
		}
	\end{center}
	
	Der Mechanismus ist keine dynamische Tensor-Glättung, sondern die Faktorisierung
	$\mathcal{C}\xrightarrow{R^*}\mathcal{F}\xrightarrow{M_{\mathrm{geom}}}24I_3$.
	Alle folgenden Resultate sind \emph{Vorbereitung} für
	Theorem~\ref{thm:tensor-restore}.
	
	\subsubsection{Vorbereitung: Definitionen und Hilfssätze}
	
	\begin{definition}[Effektiver Tensor $M_{\mathrm{eff}}$]
		\label{def:m-eff}
		Auf dem Konfigurationsraum $\mathcal{C}$ setzt man
		\[
		M_{\mathrm{eff}}(K)
		\;:=\;
		\begin{cases}
			M_{\mathrm{geom}}(I_{12},\sigma_{\mathrm{geo}}), & \text{kein Tensordefekt},\\[4pt]
			24I_3+w_pvv^{\mathsf T}, & \text{Primdefekt mit }w_p,\ \|v\|=1.
		\end{cases}
		\]
		(\texttt{M\_eff} in \texttt{EffectiveTensor.lean}).
	\end{definition}
	
	\begin{definition}[Geometrischer Tensor $M_{\mathrm{geom}}$]
		\label{def:m-geom}
		Für Vertices $V=(v_i)_{i=1}^{12}$ und Label-Zuweisung $\sigma$:
		\[
		M_{\mathrm{geom}}(V,\sigma)
		\;:=\;
		\sum_{i=1}^{12} w_{\sigma(i)}\, v_i v_i^{\mathsf T}.
		\]
	\end{definition}
	
	\begin{lemma}[Projektionssatz — effektive Ebene]
		\label{lem:tensor-projection}
		$R^*_{\mathrm{EABC}}$ löscht den Tensordefekt per Definition der Konfiguration:
		\[
		M_{\mathrm{eff}}\bigl(R^*_{\mathrm{EABC}}(K)\bigr)
		\;=\;
		M_{\mathrm{geom}}(I_{12},\sigma_{\mathrm{geo}}).
		\]
		(\texttt{M\_eff\_projection}).
	\end{lemma}
	
	\begin{theorem}[Ikosaeder-Fixpunkt — geometrische Ebene]
		\label{thm:geom-fixpoint}
		Am geometrischen Referenz-Fixpunkt gilt
		\[
		\boxed{
			M_{\mathrm{geom}}(I_{12},\sigma_{\mathrm{geo}})
			\;=\;
			24I_3
		}.
		\]
		(\texttt{geometric\_fixpoint\_isotropy}; Beweis via $R_3$-Invarianz,
		Symmetrie, $\mathrm{tr}(M)=72$ und $M_{01}=0$ aus der gewichteten Summe der
		vier $(\pm1,\pm\varphi,0)$-Ecken).
	\end{theorem}
	
	\subsubsection{Hypothesen des Hauptsatzes}
	\label{sec:main-hypotheses}
	
	Theorem~\ref{thm:tensor-restore} benötigt eine \emph{minimale} explizite
	Hypothesenliste (Lean: \texttt{prime\_norm\_full\_restoration}).
	Die folgende Tabelle trennt Definition, Eindeutigkeit und bewiesenen Inhalt.
	
	\begin{table}[h!]
		\centering
		\small
		\begin{tabular}{p{0.22\linewidth}p{0.68\linewidth}}
			\hline
			\textbf{Objekt} & \textbf{Status in Lean / Paper} \\
			\hline
			$K\in\mathcal{F}$
			& \texttt{IsEabcFixpoint}: balancierte Label-Zuweisung $\sigma_{\mathrm{geo}}$,
			keine aktive Prim-Schale, kein Tensordefekt. \\
			$p$ EABC-Prim
			& \texttt{isEabcPrime}: $p>3$ prim, $p\bmod 12\in\{1,5,7,11\}$. \\
			$w_p$
			& \textbf{Eindeutig} durch $p$: \texttt{primeWeight}\,$p=
			\texttt{familyWeight}(\texttt{classOf}\,p)\in\{1,5,7,11\}$. \\
			$v\in S^2$
			& \textbf{Nicht eindeutig}: freier Parameter in
			\texttt{extendPrimeTensor}; $\|v\|=1$ nur für $\Delta$-Aussagen
			(\texttt{delta\_prime\_defect\_at\_fixpoint}), \emph{nicht} für
			$M_{\mathrm{eff}}(R^*(K^+))=24I_3$. \\
			$R^*_{\mathrm{EABC}}$
			& \textbf{Konstruktiv}: \texttt{rStar}\,$K=\texttt{rEabc}\,K$ setzt
			$\sigma\mapsto\sigma_{\mathrm{geo}}$, löscht Prim- und Tensordefekt. \\
			$M_{\mathrm{eff}}(R^*(K^+))$
			& \textbf{Zweistufig:}
			$M_{\mathrm{eff}}(R^*(K^+))=M_{\mathrm{geom}}$ ist
			\emph{definitorisch} (\texttt{M\_eff\_projection});
			$M_{\mathrm{geom}}=24I_3$ ist \emph{bewiesener Satz}
			(\texttt{geometric\_fixpoint\_isotropy}). \\
			\hline
		\end{tabular}
		\caption{Minimale Hypothesen und epistemischer Status der zentralen Objekte.}
		\label{tab:main-hypotheses}
	\end{table}
	
	\begin{remark}[Gutachter-Checkliste]
		Die Gleichung $M_{\mathrm{eff}}(R^*(K^+))=24I_3$ ist damit \emph{kein}
		Spektralsatz über die gestörte Matrix $24I_3+w_pvv^{\mathsf T}$, sondern
		die Komposition
		\[
		\underbrace{M_{\mathrm{eff}}(R^*(K^+))=M_{\mathrm{geom}}}_{\text{Definition/Projektion}}
		\;\wedge\;
		\underbrace{M_{\mathrm{geom}}=24I_3}_{\text{Theorem (Ikosaeder)}}.
		\]
		Die Messaussagen $\Delta=w_p$ und Nicht-Isotropie sind \emph{zusätzliche},
		unabhängig bewiesene Folgerungen vor $R^*$.
	\end{remark}
	
	\subsubsection{Hauptsatz: 24-Isotropy-Restoration}
	
	\begin{theorem}[24-Isotropy-Restoration --- Hauptsatz]
		\label{thm:tensor-restore}
		Für $K\in\mathcal{F}$, EABC-Primzahl $p$ und beliebiges $v\in\mathbb{R}^3$ mit
		$K^+=\mathrm{Extend}(p,v)\,K$ gilt die zentrale Restauration
		(unabhängig von $\|v\|$; siehe Abschnitt~\ref{sec:what-we-prove})
		\[
		\boxed{
			M_{\mathrm{eff}}\bigl(R^*_{\mathrm{EABC}}(K^+)\bigr) = 24I_3
		}.
		\]
		Vollständige Satzkette (Lean: Tabelle~\ref{tab:phase-a-chain}):
		\begin{align*}
			M_{\mathrm{geom}}(I_{12},\sigma_{\mathrm{geo}}) &= 24I_3,
			&
			M_{\mathrm{eff}}(K^+) &= 24I_3+w_pvv^{\mathsf T},
			\\
			\Delta\bigl(M_{\mathrm{eff}}(K^+)\bigr) &= w_p,
			&
			w_p>0 &\Rightarrow M_{\mathrm{eff}}(K^+)\ \text{nicht isotrop},
			\\
			R^*_{\mathrm{EABC}}(K^+) &\in \mathcal{F},
			&
			M_{\mathrm{eff}}\bigl(R^*_{\mathrm{EABC}}(K^+)\bigr) &= 24I_3.
		\end{align*}
		(\texttt{prime\_norm\_full\_restoration} in \texttt{TensorRestoration.lean}).
	\end{theorem}
	
	\begin{proof}
		Theorem~\ref{thm:prime-delta} und \texttt{delta\_prime\_defect\_at\_fixpoint}
		(vor $R^*$),
		\texttt{rankOne\_defect\_not\_isotropic} (Nicht-Isotropie),
		Theorem~\ref{thm:prime-norm} bzw.\ \texttt{prime\_norm\_full\_restoration}
		(Label und Tensor nach $R^*$),
		Lemma~\ref{lem:tensor-projection} (definitorisch, Ebene~1),
		Theorem~\ref{thm:geom-fixpoint} (geometrisch, Ebene~2).
	\end{proof}
	
	\begin{corollary}[Globale Peano-Fortsetzung]
		\label{cor:shellstack-global}
		Auf einem Schalenstapel $\mathcal{S}$ mit lokalen Fixpunkten gilt nach komponentenweisem
		$R^*$ schalenweise $M_{\mathrm{eff}}=24I_3$
		(\texttt{all\_shells\_tensor\_restored}).
	\end{corollary}
	
	\begin{remark}[Gutachter-Einordnung]
		$M_{\mathrm{eff}}(R^*(K))=24I_3$ ist \emph{kein} mystischer Spektralsatz:
		die Projektion ersetzt die defekte Konfiguration durch den geometrischen Fixpunkt;
		dessen Isotropie ist der inhaltliche Satz (Theorem~\ref{thm:geom-fixpoint}).
		$P_I$ (Voronoi-Pullback gestörter Vertices) ist eine geometrische Realisierung
		von $R^*$, nicht der konzeptionelle Kern.
	\end{remark}
	
	\begin{remark}[$P_I$ als geometrische Realisierung (optional)]
		\label{rem:pi-optional}
		Für gestörte Richtungen $v\in S^2$ setzt
		$P_I(v):=\arg\max_{r_j\in I_{12}}\langle v,r_j\rangle$.
		Bei hinreichend kleiner Kuss-Störung zieht $P_I$ die Vertex-Menge auf $I_{12}$
		zurück; zusammen mit labelblindem $R_{\mathrm{EABC}}$ erhält man dieselbe
		Projektionssemantik wie oben.
	\end{remark}
	
	\begin{lemma}[Multiketten-Imbalance nach dem Primschritt]
		\label{lem:prime-imbalance}
		Aus $K_n\in\mathcal{F}$ und $p=P(2n)$ folgt für die Familienzählung mit virtueller
		Primkugel $(4,3,3,3)$ statt $(3,3,3,3)$.
	\end{lemma}
	
	%=============================================================================
	\section{Was dieses Manuskript beweist --- und was nicht}
	\label{sec:what-we-prove}
	%=============================================================================
	
	Diese Seite fasst den epistemischen Stand für externe Leser zusammen.
	Details: Hypothesen (Tabelle~\ref{tab:main-hypotheses}),
	Abhängigkeiten (Abschnitt~\ref{sec:formal-deps}),
	Hauptsatz (Theorem~\ref{thm:tensor-restore}).
	
	\subsection{Objektrollen}
	
	\begin{table}[h!]
		\centering
		\small
		\begin{tabular}{p{0.18\linewidth}p{0.74\linewidth}}
			\hline
			\textbf{Objekt} & \textbf{Rolle} \\
			\hline
			$24I_3$ & geometrischer Fixpunkt (Referenz-Isotropie) \\
			$w_p$ & Primlabel, eindeutig aus $p\bmod 12$ \\
			$v$ & freier Richtungsparameter (nicht durch $p$ festgelegt) \\
			$vv^{\mathsf T}$ & Rang-1-Defektkomponente in $M_{\mathrm{eff}}(K^+)$ \\
			$R^*_{\mathrm{EABC}}$ & Retraktionsoperator (projektive Konfigurationsabbildung) \\
			$\Delta$ & Messgröße der Anisotropie (Ordnungsparameter, $\Delta=w_p$ vor $R^*$) \\
			\hline
		\end{tabular}
		\caption{Begriffslexikon: Rolle der zentralen Objekte.}
		\label{tab:object-roles}
	\end{table}
	
	\subsection{Tatsächliche logische Struktur}
	
	\textbf{Hauptpfad (Restauration):} unabhängig von $v$ und unabhängig von $\Delta$.
	\[
	\texttt{IsEabcFixpoint}\,K
	\;\wedge\;
	\texttt{isEabcPrime}\,p
	\;\Longrightarrow\;
	\texttt{prime\_norm\_full\_restoration}
	\;\Longrightarrow\;
	M_{\mathrm{eff}}(R^*(K^+))=24I_3.
	\]
	Der Tensor-Teil faktorisiert als
	\[
	M_{\mathrm{eff}}(R^*(K^+))=M_{\mathrm{geom}}
	\quad\text{(Definition)}
	\qquad\text{und}\qquad
	M_{\mathrm{geom}}=24I_3
	\quad\text{(Theorem)}.
	\]
	\textbf{Keine neue Information} steckt in $M_{\mathrm{eff}}(R^*(K^+))=M_{\mathrm{geom}}$;
	der \emph{mathematische Gehalt} des Hauptsatzes liegt vollständig in
	$M_{\mathrm{geom}}(I_{12},\sigma_{\mathrm{geo}})=24I_3$
	(\texttt{geometric\_fixpoint\_isotropy}).
	
	\medskip\noindent
	\textbf{Parallelpfad (Messung):} unabhängig vom Restaurationsbeweis.
	\[
	K^+ \;\longrightarrow\; M_{\mathrm{eff}}(K^+)=24I_3+w_pvv^{\mathsf T}
	\;\longrightarrow\; \Delta(M_{\mathrm{eff}}(K^+))=w_p
	\;\longrightarrow\; \text{Messgröße}.
	\]
	Die Lemmas \texttt{delta\_prime\_defect\_at\_fixpoint} und
	\texttt{rankOne\_defect\_not\_isotropic} gehören hierher; sie sind
	\emph{Vorbereitung/Messung}, keine logische Voraussetzung des Gleichheitsbeweises.
	
	\paragraph{Konsequenz: $v$-Unabhängigkeit.}
	Da $\|v\|$ für \texttt{prime\_norm\_full\_restoration} nicht benötigt wird,
	entfernt $R^*$ die gesamte Klasse der zulässigen Rang-1-Erweiterungen
	$24I_3+w_pvv^{\mathsf T}$ unabhängig von der konkreten Richtung $v$.
	Der Hauptsatz ist damit \emph{stärker} als eine rein vektor-spezifische Aussage.
	
	\subsection{Drei epistemische Ebenen}
	
	\begin{table}[h!]
		\centering
		\small
		\begin{tabular}{p{0.16\linewidth}p{0.78\linewidth}}
			\hline
			\textbf{Ebene} & \textbf{Inhalt} \\
			\hline
			1 --- Definitionen &
			$R^*$, $M_{\mathrm{eff}}$, $w_p$, $v$, \texttt{extendPrimeTensor},
			\texttt{M\_eff\_projection}. \\
			2 --- Bewiesene Struktur &
			$M_{\mathrm{geom}}=24I_3$;
			$\Delta=w_p$ (Messung);
			$M_n^{\mathrm{sep}}=4^n$ und Skalierungsinvarianz (Phase~B). \\
			3 --- Offene Programme &
			$\dim_{\mathrm{box}}$; Ball-Überdeckung; Grenzwert $n\to\infty$;
			Monopol-Physik ($Q_m$, Dirac--Schwinger);
			Phase-C-Hypothesen (QCD-Lesart, superkritische Normschale, AdS/CFT-Bulk). \\
			\multicolumn{2}{@{}l@{}}{\emph{Teilweise geschlossen (Phase~C):}
				$\iota_n$ für $n\le 3$ (\texttt{cardinalShellEmbedding\_one/two/three}),
				Ikosaeder-Separation $\varepsilon=1$ (\texttt{IcosahedronEmbedding.lean}),
				vier Mathlib-U(1)-Anker (\texttt{DiracMonopole.lean}).} \\
			\hline
		\end{tabular}
		\caption{Drei Ebenen: Definition, Beweis, offenes Programm.}
		\label{tab:three-epistemic-layers}
	\end{table}
	
	\subsection{Was bewiesen ist}
	
	\begin{itemize}
		\item \textbf{Geometrischer Kern:}
		$M_{\mathrm{geom}}(I_{12},\sigma_{\mathrm{geo}})=24I_3$
		(\texttt{geometric\_fixpoint\_isotropy}).
		\item \textbf{Projektive Semantik:}
		$M_{\mathrm{eff}}(R^*(K))=M_{\mathrm{geom}}$ per Definition der Konfiguration
		(\texttt{M\_eff\_projection}).
		\item \textbf{24-Isotropy-Restoration:}
		unter \texttt{IsEabcFixpoint} und \texttt{isEabcPrime} gilt
		$M_{\mathrm{eff}}(R^*(K^+))=24I_3$ für \emph{jedes} $v$
		(\texttt{prime\_norm\_full\_restoration}).
		\item \textbf{Messung:} $\Delta(M_{\mathrm{eff}}(K^+))=w_p$ bei $\|v\|=1$.
		\item \textbf{Peano:} schalenweise Restauration (\texttt{all\_shells\_tensor\_restored}).
		\item \textbf{Phase~B (Nebenzweig):}
		$M_n^{\mathrm{sep}}=C_n=4^n$;
		$\ln M_n^{\mathrm{sep}}/\ln r_n=\ln4/\ln\phi$ (Skalierungsinvarianzsatz, keine Box-Dimension).
		\item \textbf{Phase~C (Teilstand):}
		metrische Einbettungen $\iota_n$ für $n\in\{1,2,3\}$ mit
		$\varepsilon_n\in\{1,\varphi^{-2},\varphi^{-3}\}$
		(\texttt{cardinalShellEmbedding\_one/two/three});
		Ikosaeder-Repräsentanten-Separation $\varepsilon=1$
		(\texttt{icosahedronShellEmbedding\_one\_sep});
		vier Mathlib-U(1)-Anker im Dirac-Monopol-Scaffold
		(\texttt{DiracMonopole.lean}, Abschnitt~\ref{sec:dirac-monopole}).
	\end{itemize}
	
	\subsection{Was nicht bewiesen ist}
	
	\begin{itemize}
		\item Dass $M_{\mathrm{eff}}(R^*(K^+))=24I_3$ ein Spektralsatz über die
		gestörte Matrix $24I_3+w_pvv^{\mathsf T}$ ist (es ist Projektion + Fixpunkt).
		\item Metrische Minimalüberdeckung $M_n(\varepsilon_n)$ und Minkowski--Bouligand-Dimension.
		\item Vollständige Einbettung $\iota_n$ für alle $n$ und Grenzwert $n\to\infty$
		(für $n\le 3$ siehe Phase~C-Teilstand oben).
		\item Dirac-Ladungsquantisierung $Q_m$, Spin-Eis-Monopole, AdS/CFT-Bulk-Formalismus
		(nur Hypothesen/Remark, Abschnitt~\ref{sec:dirac-monopole}).
		\item Physikalische Deutung von $w_p$, $R^*$ oder der Peano-Schachtelung
		(einschließlich Hypothesen~\ref{hyp:qcd-gauge-fix} und~\ref{hyp:shell-collapse}).
		\item Dass $p_{\mathrm{krit}}=173$ oder $Z_{\mathrm{krit}}\approx 173$ aus der
		EABC-Arithmetik folgt.
		\item Konvergenz unendlicher Residuensummen aus $D_{EABC}<3$ allein.
	\end{itemize}
	
	%=============================================================================
	\section{Formale Verifikation}
	\label{sec:formal-verification}
	%=============================================================================
	
	Die zentrale Satzkette ist in \texttt{eabc-renorm} (Lean~4.29, Mathlib~4.29)
	maschinell verifiziert; \texttt{lake build} ist grün, \texttt{sorry}\,$=\,0$.
	Hauptsatz: Theorem~\ref{thm:tensor-restore} (24-Isotropy-Restoration).
	Überblick: Abschnitt~\ref{sec:what-we-prove}.
	
	\begin{center}
		\small
		\begin{minipage}[t]{0.30\linewidth}
			\fbox{\parbox{\linewidth}{
					\textbf{PHASE A}\\[4pt]
					$\checkmark$ vollständig formalisiert\\
					$\checkmark$ Lean-belegt\\
					$\checkmark$ \textbf{zentrale Aussage}\\
					\textit{Hauptsatz~\ref{thm:tensor-restore}}
			}}
		\end{minipage}
		\hfill
		\begin{minipage}[t]{0.30\linewidth}
			\fbox{\parbox{\linewidth}{
					\textbf{PHASE B}\\[4pt]
					$\checkmark$ kombinatorische Struktur\\
					$\checkmark$ Skalierungsinvarianz\\
					$\times$ keine Box-Dimension\\
					\textit{Nebenzweig / Einordnung}
			}}
		\end{minipage}
		\hfill
		\begin{minipage}[t]{0.30\linewidth}
			\fbox{\parbox{\linewidth}{
					\textbf{PHASE C}\\[4pt]
					$\checkmark$ $\iota_n$, $n\le 3$\\
					$\checkmark$ Ikosaeder-Separation\\
					$\checkmark$ Dirac-Anker (Mathlib)\\
					$\circ$ Box-Dim / $n\to\infty$\\
					\textit{Teilstand + Forschungsprogramm}
			}}
		\end{minipage}
	\end{center}
	
	\subsection{Drei-Phasen-Architektur}
	\label{sec:three-phases}
	
	Die Verifikation ist bewusst in drei Schichten gegliedert.
	Phase~A und Phase~B sind als geschlossener Stand dokumentiert;
	Phase~C hat einen dokumentierten Teilstand ($n\le 3$, Ikosaeder, Dirac-Scaffold),
	Ball-Überdeckung und Grenzwert $n\to\infty$ bleiben offen.
	
	\paragraph{Phase~A (geschlossen): geometrischer Fixpunkt und Tensor-Renormierung.}
	Kernsubstanz der physikalisch-geometrischen Lesart:
	\[
	M_{\mathrm{geom}}(I_{12},\sigma_{\mathrm{geo}})=24I_3,
	\qquad
	\Delta(M^+)=w_p,
	\qquad
	w_p>0 \Rightarrow M^+\ \text{nicht isotrop},
	\qquad
	R^*(K^+)\in\mathcal{F},\ M_{\mathrm{eff}}=24I_3.
	\]
	Peano-Multi-Schale: \texttt{ShellStack}, \texttt{all\_shells\_tensor\_restored}.
	Hauptsatz: Theorem~\ref{thm:tensor-restore}; Details in
	Abschnitt~\ref{sec:phase-a-chain} und Tabelle~\ref{tab:two-levels}.
	
	\paragraph{Phase~B (Nebenzweig, eingefroren): diskrete fraktale Zustandsstruktur.}
	\emph{Einordnung, nicht Hauptsatz.}
	Kombinatorischer Wortbaum, kanonische und minimale \emph{separierte} Überdeckung:
	\[
	C_n=4^n,
	\qquad
	M_n^{\mathrm{sep}}=C_n=4^n,
	\qquad
	\frac{\ln M_n^{\mathrm{sep}}}{\ln r_n}
	\;=\;
	\frac{\ln 4}{\ln\phi}.
	\]
	Das ist ein formaler Minimalcover-Satz auf endlichen Schalen, kein
	$\mathbb{R}^3$-Infimum und keine Box-Dimension.
	
	\paragraph{Phase~C (Teilstand + Anschlussarbeit).}
	Metrische Realisierung erfordert eine Einbettung
	$\iota_n:\texttt{ShellVertex}(n)\to\mathbb{R}^3$ und eine Separationsbedingung
	\[
	u\neq v
	\;\Longrightarrow\;
	\|\iota_n(u)-\iota_n(v)\|\ge\varepsilon_n.
	\]
	Für $n\in\{1,2,3\}$ sind \texttt{cardinalShellEmbedding\_one/two/three}
	und die Ikosaeder-Repräsentanten-Einbettung
	\texttt{icosahedronShellEmbedding\_one} mit $\varepsilon=1$ bewiesen
	(\texttt{IcosahedronEmbedding.lean}).
	Erst für alle $n$ und mit Ball-Überdeckungen folgt die metrische Version
	$M_n(\varepsilon_n)=C_n=4^n$ und ein Grenzwert
	$\lim_{n\to\infty}\log M_n/\log r_n$ im Sinne der Minkowski--Bouligand-Dimension.
	\texttt{DiracMonopole.lean} liefert vier Mathlib-U(1)-Anker als Scaffold;
	\texttt{GeometryScaffold.lean} liefert einen Mathlib-Ptolemy-Anker und vier
	klassische Geometrie-Hypothesen (Abschnitt~\ref{sec:geometry-scaffold}).
	Monopol-Physik bleibt Hypothese (Abschnitt~\ref{sec:dirac-monopole}).
	Flyspeck-Analogie: \texttt{VerificationCertificate.lean} trennt Lean-Beweise
	von externen Python-Zertifikaten (\texttt{eabc\_icosahedron\_test.py}) ---
	die Zertifikate ersetzen keine Lean-Theoreme
	(Abschnitt~\ref{sec:flyspeck-certificates}).
	
	\subsection{Phase~A: Kernsatz und Lean-Satzkette}
	\label{sec:phase-a-chain}
	
	Die Phase-A-Verifikation faktorisiert strikt in \emph{Störung},
	\emph{Messung}, \emph{Projektion} und \emph{geometrischen Fixpunkt}:
	
	\begin{center}
		\fbox{%
			\begin{minipage}{0.92\linewidth}
				\vspace{2mm}
				\textbf{Theorem~\ref{thm:tensor-restore} (24-Isotropy-Restoration).}
				\vspace{2mm}
				\[
				\underbrace{M_{\mathrm{eff}}(K^+)=24I_3+w_pvv^{\mathsf T}}_{\text{Störung}}
				\;\Longrightarrow\;
				\underbrace{\Delta(M_{\mathrm{eff}}(K^+))=w_p}_{\text{messbar}}
				\;\xrightarrow{\;R^*\;}
				\underbrace{M_{\mathrm{eff}}(R^*(K^+))=24I_3}_{\text{restauriert}},
				\]
				\[
				\text{weil}\quad
				\underbrace{M_{\mathrm{geom}}(I_{12},\sigma_{\mathrm{geo}})=24I_3}_{\text{geometrischer Fixpunkt}}.
				\]
				\vspace{2mm}
			\end{minipage}%
		}
	\end{center}
	
	\begin{table}[h!]
		\centering
		\small
		\begin{tabular}{cll}
			\hline
			\textbf{Schritt} & \textbf{Aussage} & \textbf{Lean} \\
			\hline
			1 & $M_{\mathrm{geom}}(I_{12},\sigma_{\mathrm{geo}})=24I_3$ &
			\texttt{geometric\_fixpoint\_isotropy} \\
			2 & $M_{\mathrm{eff}}(K^+)=24I_3+w_pvv^{\mathsf T}$ &
			\texttt{M\_eff\_with\_defect}, \texttt{extendPrimeTensor} \\
			3 & $\Delta(M_{\mathrm{eff}}(K^+))=w_p$; $w_p>0\Rightarrow$ nicht isotrop &
			\texttt{delta\_prime\_defect\_at\_fixpoint}, \texttt{rankOne\_defect\_not\_isotropic} \\
			4 & $M_{\mathrm{eff}}(R^*(K))=M_{\mathrm{geom}}$ (Projektion) &
			\texttt{M\_eff\_projection} \\
			5 & $R^*(K^+)\in\mathcal{F}$; $M_{\mathrm{eff}}(R^*(K^+))=24I_3$ &
			\texttt{prime\_norm\_full\_restoration} \\
			6 & Peano-global: schalenweise $24I_3$ &
			\texttt{all\_shells\_tensor\_restored} \\
			\hline
		\end{tabular}
		\caption{Phase-A-Satzkette: geometrischer Fixpunkt, effektive Störung, Projektion, Restauration.}
		\label{tab:phase-a-chain}
	\end{table}
	
	\paragraph{Gutachter-Hinweis.}
	$M_{\mathrm{eff}}(R^*(K))=24I_3$ ist \emph{kein} Spektralsatz über die gestörte Matrix,
	sondern die Komposition aus definitorischer Projektion (Schritt~4) und dem
	inhaltlichen Ikosaeder-Fixpunkt (Schritt~1).
	$P_I$ (Voronoi-Pullback) ist optional und nicht Teil der Lean-Verifikation.
	
	\subsection{Formale Abhängigkeiten}
	\label{sec:formal-deps}
	
	Gerichteter Abhängigkeitsgraph der Phase-A-Verifikation.
	\textbf{Def} = definitorisch/axiombasiert in Lean;
	\textbf{Thm} = bewiesenes Lemma/Theorem.
	Hypothesen des Hauptsatzes: Tabelle~\ref{tab:main-hypotheses}.
	
	\begin{center}
		\small
		\begin{tabular}{c}
			\textit{Def:} \texttt{rStar}, \texttt{extendPrimeTensor}, \texttt{M\_eff}, \texttt{primeWeight} \\[6pt]
			$\downarrow$ \\[6pt]
			\textbf{Thm:} \texttt{geometric\_fixpoint\_isotropy} \quad ($M_{\mathrm{geom}}=24I_3$) \\[6pt]
			$\downarrow$ \\[6pt]
			\textbf{Def:} \texttt{M\_eff\_projection} \quad ($M_{\mathrm{eff}}(R^*K)=M_{\mathrm{geom}}$) \\[6pt]
			$\swarrow$ \hspace{2cm} $\searrow$ \\[4pt]
			\begin{minipage}{0.42\linewidth}
				\raggedright
				\textit{parallel (Messung, vor $R^*$):}\\
				\textbf{Thm:} \texttt{delta\_prime\_defect\_at\_fixpoint}\\
				$\leftarrow$ \texttt{primeWeight\_pos}\\
				$\leftarrow$ \texttt{rankOne\_defect\_not\_isotropic}
			\end{minipage}
			\hspace{0.04\linewidth}
			\begin{minipage}{0.42\linewidth}
				\raggedright
				\textbf{Thm:} \texttt{M\_eff\_after\_rStar}\\
				($M_{\mathrm{eff}}(R^*K^+)=24I_3$)
			\end{minipage}
			\\[8pt]
			$\downarrow$ \\[6pt]
			\textbf{Thm:} \texttt{prime\_norm\_restoration} \quad (Label: $R^*(K^+)\in\mathcal{F}$) \\[6pt]
			$\downarrow$ \\[6pt]
			\fbox{\textbf{Thm:} \texttt{prime\_norm\_full\_restoration} \; (Hauptsatz)} \\[6pt]
			$\downarrow$ \\[6pt]
			\textbf{Thm:} \texttt{all\_shells\_tensor\_restored} \quad (Peano-global)
		\end{tabular}
	\end{center}
	
	\begin{table}[h!]
		\centering
		\small
		\begin{tabular}{lll}
			\hline
			\textbf{Status} & \textbf{Inhalt} & \textbf{Lean / Anmerkung} \\
			\hline
			$\checkmark$ bewiesen & 24-Isotropy-Restoration & \texttt{prime\_norm\_full\_restoration} \\
			$\checkmark$ bewiesen & Peano-ShellStack-Induktion & \texttt{all\_shells\_tensor\_restored} \\
			$\checkmark$ bewiesen & kombinatorische Skalierung (Phase~B) &
			\texttt{fractal\_ratio\_via\_minimal\_separated\_cover} \\
			$\checkmark$ bewiesen & $\mathbb{R}^3$-Einbettung $n\le 3$ &
			\texttt{cardinalShellEmbedding\_one/two/three} \\
			$\checkmark$ bewiesen & Ikosaeder-Separation $\varepsilon=1$ &
			\texttt{icosahedronShellEmbedding\_one\_sep} \\
			$\checkmark$ bewiesen & Dirac-Monopol Mathlib-Anker (4) &
			\texttt{DiracMonopole.milestone\_green\_*} \\
			$\circ$ offen & metrische Box-Dimension & Phase~C \\
			$\circ$ offen & Ball-Überdeckung, Grenzwert $n\to\infty$ & Phase~C \\
			$\circ$ offen & Instabilitäts-API / \texttt{ShellInstability} &
			\texttt{MetricSeparationLossExists} offen \\
			$\circ$ offen & Monopol-Physik ($Q_m$, Dirac--Schwinger) &
			\texttt{DiracMonopole} (3 Hypothesen) \\
			$\circ$ offen & QCD-/Stella-Octangula-Lesart & Hypothese~\ref{hyp:qcd-gauge-fix} \\
			$\circ$ offen & superkritische Normschale & Hypothese~\ref{hyp:shell-collapse} \\
			\hline
		\end{tabular}
		\caption{Epistemischer Status: was bewiesen ist und was bewusst offen bleibt.}
		\label{tab:epistemic-status}
	\end{table}
	
	\begin{table}[h!]
		\centering
		\small
		\begin{tabular}{p{0.38\linewidth}p{0.30\linewidth}p{0.24\linewidth}}
			\hline
			\textbf{Mathematische Aussage} & \textbf{Lean-Theorem} & \textbf{Datei} \\
			\hline
			$M_{\mathrm{geom}}(I_{12},\sigma_{\mathrm{geo}})=24I_3$
			& \texttt{geometric\_fixpoint\_isotropy} & \texttt{IcosahedronSymmetry.lean} \\
			$M^+=24I_3+w_pvv^{\mathsf T}$
			& \texttt{rankOnePerturbation} & \texttt{PrimeDefect.lean} \\
			$\Delta(M^+)=w_p$
			& \texttt{delta\_prime\_defect\_at\_fixpoint} & \texttt{PrimeDefect.lean} \\
			$w_p>0\Rightarrow M^+$ nicht isotrop
			& \texttt{rankOne\_defect\_not\_isotropic} & \texttt{TensorRestoration.lean} \\
			$R^*(K^+)\in\mathcal{F}$, $M_{\mathrm{eff}}=24I_3$
			& \texttt{prime\_norm\_full\_restoration} & \texttt{TensorRestoration.lean} \\
			$M_{\mathrm{eff}}(R^*(K))=M_{\mathrm{geom}}(\mathrm{ref})$
			& \texttt{M\_eff\_projection} & \texttt{EffectiveTensor.lean} \\
			\hline
			\multicolumn{3}{l}{\textbf{Peano-Erweiterung (Multi-Schale, Phase~A)}} \\
			\hline
			Schalenstapel $\mathcal{S}$
			& \texttt{structure ShellStack} & \texttt{ShellStack.lean} \\
			Lokaler Schalen-Retrakt
			& \texttt{shell\_defect\_restoration} & \texttt{ShellStack.lean} \\
			Globaler Projektionssatz
			& \texttt{all\_shells\_restored} & \texttt{ShellStack.lean} \\
			Globale Tensor-Restauration
			& \texttt{all\_shells\_tensor\_restored} & \texttt{ShellStack.lean} \\
			\hline
			\multicolumn{3}{l}{\textbf{Phase B: Wortbaum und Skalierungsexponent}} \\
			\hline
			$C_n=\texttt{card}(\texttt{ShellVertex}\,n)=4^n$
			& \texttt{canonicalCoveringNumber\_eq\_four\_pow} & \texttt{ShellGraph.lean} \\
			Kanonische Überdeckung existiert ($C_n$)
			& \texttt{canonical\_cover\_exists} & \texttt{ShellSeparation.lean} \\
			$\log C_n/\log r_n=\log4/\log\phi$
			& \texttt{fractal\_ratio\_via\_canonical\_cover} & \texttt{FractalDimension.lean} \\
			$M_n^{\mathrm{sep}}=C_n=4^n$ (separierte Zellen)
			& \texttt{minimal\_separated\_cover\_eq\_four\_pow} & \texttt{ShellMinimalCover.lean} \\
			$\#V_n\le\#\text{cells}$ bei Separation
			& \texttt{separated\_cover\_lower\_bound} & \texttt{ShellMinimalCover.lean} \\
			$M_n^{\mathrm{sep}}=C_n$ (Abstraktion)
			& \texttt{separation\_implies\_minimal\_cover} & \texttt{ShellSeparation.lean} \\
			$\log M_n^{\mathrm{sep}}/\log r_n=\log4/\log\phi$
			& \texttt{fractal\_ratio\_via\_minimal\_separated\_cover} & \texttt{FractalDimension.lean} \\
			$Z(n)=4^n$, $r_n=\phi^n$ (Modell)
			& \texttt{eabcZustand}, \texttt{metrischeSkalierung} & \texttt{FractalDimension.lean} \\
			$\log(4^n)/\log(\phi^n)=\log4/\log\phi$
			& \texttt{fractal\_dimension\_invariant} & \texttt{FractalDimension.lean} \\
			$\ln4/\ln\phi<3$ für $\phi=(1+\sqrt5)/2$
			& \texttt{fractal\_dimension\_lt\_three} & \texttt{FractalDimension.lean} \\
			\hline
			\multicolumn{3}{l}{\textbf{Phase C: metrische Realisierung (Teilstand)}} \\
			\hline
			$\iota_1$ mit Separation $\varepsilon=1$
			& \texttt{cardinalShellEmbedding\_one} & \texttt{ShellEmbedding.lean} \\
			$\iota_2$, $\varepsilon=\varphi^{-2}$
			& \texttt{cardinalShellEmbedding\_two} & \texttt{ShellEmbedding.lean} \\
			$\iota_3$, $\varepsilon=\varphi^{-3}$
			& \texttt{cardinalShellEmbedding\_three} & \texttt{ShellEmbedding.lean} \\
			Ikosaeder-Repräsentanten, $\varepsilon=1$
			& \texttt{icosahedronShellEmbedding\_one\_sep} & \texttt{IcosahedronEmbedding.lean} \\
			$\varphi$-Skala erreichbar ($n\le 3$)
			& \texttt{phiScaleAchievable\_one/two/three} & \texttt{ShellSeparationScale.lean} \\
			$\neg\,\mathrm{ShellSeparationLoss}(n)$ für $n\le 3$
			& \texttt{milestone\_green\_no\_loss\_one/two/three} & \texttt{PhaseCBlueprint.lean} \\
			Ikosaeder-Meilenstein (Blueprint)
			& \texttt{milestone\_green\_icosahedron} & \texttt{PhaseCBlueprint.lean} \\
			U(1)-Periodizität (Mathlib-Anker)
			& \texttt{milestone\_green\_u1\_period} & \texttt{DiracMonopole.lean} \\
			$S^1\simeq\mathrm{AddCircle}(2\pi)$ (Mathlib-Anker)
			& \texttt{milestone\_green\_circle\_homeomorph} & \texttt{DiracMonopole.lean} \\
			Schleifen-Homotopie (Mathlib-Anker)
			& \texttt{milestone\_green\_loop\_homotopy} & \texttt{DiracMonopole.lean} \\
			$\pi_1$-Vorbereitung (Mathlib-Anker)
			& \texttt{milestone\_green\_fundamental\_group} & \texttt{DiracMonopole.lean} \\
			Ptolemy (Mathlib-Anker)
			& \texttt{milestone\_green\_mathlib\_ptolemy} & \texttt{GeometryScaffold.lean} \\
			Geometrie-Hypothesen (4$\times$, rot)
			& \texttt{MorleyEquilateralHypothesis}, \ldots &
			\texttt{GeometryScaffold.lean} \\
			Flyspeck-Metadaten (Phase~C)
			& \texttt{phaseCMilestones} & \texttt{VerificationCertificate.lean} \\
			\hline
		\end{tabular}
		\caption{Verifizierte Satzkette: Phase~A (Tensor/Peano), Phase~B (kombinatorische Fraktalstruktur) und Phase~C (Teilstand).}
		\label{tab:lean-theorems}
	\end{table}
	
	\subsection{Peano-Schachtelung und globaler Retrakt}
	
	Eine einzelne Schale ist in Lean abgeschlossen.
	Für geschachtelte Normschalen wird der Konfigurationsraum zu einem
	\emph{indexbasierten Stapel} $\mathcal{S}=\{K^{(i)}\}_{i<n}$ über
	$\mathbb{N}$ (Peano-Rekursion) erweitert.
	Der globale Operator $\mathcal{R}^*$ wirkt komponentenweise:
	auf markierten Schalen $i$ mit aktivem Primdefekt gilt
	\[
	\mathcal{R}^*(K^{(i)}) \;=\; R^*_{\mathrm{EABC}}\!\bigl(\mathrm{Extend}(p)\,K^{(i)}\bigr),
	\]
	auf unberührten Schalen bleibt $K^{(i)}$ unverändert.
	Da jede Schale den lokalen Satz \texttt{shell\_defect\_restoration} erfüllt,
	folgt der globale Satz \texttt{all\_shells\_restored} ohne zusätzliche
	Analyse-Hypothesen.
	
	\subsection{Metrische Interpretation (Phase~B, Nebenzweig)}
	
	Die kombinatorische Verifikation hängt \emph{nicht} von metrischen Radien ab.
	Für physikalische Lesarten lässt sich dennoch eine rein interpretative
	Skalierungsfunktion einführen, z.\,B.
	\[
	r_n \;=\; \phi^n\,r_0,
	\qquad \phi=\tfrac{1+\sqrt5}{2},
	\]
	ohne den algebraischen Restaurationssatz zu verändern:
	\texttt{all\_shells\_restored} bleibt exakt, weil $R^*$ auf der
	Label-/Defektebene operiert und $M_{\mathrm{geom}}=24I_3$ metrikunabhängig ist.
	
	\subsection{Skalierungsexponent $D_{EABC}$ (Phase~B, Nebenzweig; keine Box-Dimension)}
	\label{sec:fractal-dimension}
	
	\paragraph{Was mathematisch sauber ist.}
	Unter der \emph{Modellannahme}
	\[
	Z(n)=4^n,\qquad r_n=\phi^n\qquad(\phi>1)
	\]
	folgt das exakte Skalierungsgesetz
	\[
	\frac{\ln Z(n)}{\ln r_n}
	\;=\;
	\frac{n\ln 4}{n\ln\phi}
	\;=\;
	\frac{\ln 4}{\ln\phi}.
	\]
	Lean beweist genau dies in \texttt{FractalDimension.lean}:
	\texttt{fractal\_dimension\_invariant} liefert für jedes $n>0$
	\[
	\frac{\log(4^n)}{\log(\phi^n)}
	\;=\;
	\frac{\log 4}{\log\phi}.
	\]
	Das ist ein \textbf{Skalierungsinvarianzsatz}, kein Grenzwertsatz
	$\epsilon\to 0$ der klassischen Minkowski--Bouligand-Dimension
	\[
	D_B=\lim_{\epsilon\to0}\frac{\log M(\epsilon)}{\log(1/\epsilon)}.
	\]
	
	\paragraph{Abgrenzung zu Phase~C.}
	Phase~B beweist $M_n^{\mathrm{sep}}=C_n=4^n$ und damit
	$\ln M_n^{\mathrm{sep}}/\ln r_n=\ln4/\ln\phi$ als kombinatorisches
	Skalierungsgesetz.
	Daraus folgt \emph{nicht automatisch}, dass
	$D_{EABC}=\ln4/\ln\phi$ die Box-Dimension eines in $\mathbb{R}^3$
	eingebetteten Objekts ist (Phase~C).
	Dafür bräuchte man $\iota_n:\texttt{ShellVertex}(n)\to\mathbb{R}^3$,
	metrische Separations- und Überdeckungsbedingungen sowie einen Grenzwert
	$\epsilon\to 0$.
	
	\paragraph{Numerische Kennzahl und Sierpiński-Vergleich.}
	Für $\phi=(1+\sqrt5)/2$ gilt
	\[
	D_{EABC}:=\frac{\ln 4}{\ln\phi}\approx 2{,}88137,
	\qquad
	3-D_{EABC}\approx 0{,}12.
	\]
	Lean beweist algebraisch $D_{EABC}<3$ (\texttt{fractal\_dimension\_lt\_three},
	über $\phi^3=2+\sqrt5>4$).
	Zum Vergleich: beim Sierpiński-Tetraeder ist
	$D_S=\ln4/\ln2=2$ (vier Kopien, Maßstab $1/2$).
	Da $\phi<2$, liegt $D_{EABC}$ \emph{über} $D_S$: das EABC-Modell ist
	kombinatorisch dichter als das klassische Sierpiński-Beispiel --- anschaulich
	näher an raumfüllend, aber formal noch nicht als geometrisches Volumenargument
	bewiesen.
	
	\paragraph{Konvergenz der Residuen (Vorsicht).}
	Die Bedingung $D_{EABC}<3$ impliziert \emph{nicht} automatisch, dass unendliche
	Summen lokaler Krümmungs- oder Defektresiduen konvergieren.
	Dafür bräuchte man zusätzlich ein Abfallgesetz $R_n\sim r_n^{-\alpha}$ und die
	Analyse von $\sum_n R_n$.
	Für das Paper gilt daher die robustere Formulierung:
	\begin{quote}
		$D_{EABC}<3$ liefert eine \emph{notwendige geometrische Voraussetzung} für die
		Endlichkeit globaler Retraktionsprozesse; die tatsächliche Konvergenz hängt
		zusätzlich vom Abfallgesetz der lokalen Defektgewichte ab.
	\end{quote}
	Die algebraische Retraktionskette (\texttt{all\_shells\_tensor\_restored},
	\texttt{prime\_norm\_full\_restoration}) bleibt davon unabhängig und ist
	bereits auf jeder endlichen Schale geschlossen.
	
	\paragraph{Schritt 1 (geschlossen): kanonische Zellenüberdeckung.}
	In \texttt{ShellGraph.lean}:
	\[
	\texttt{ShellVertex}(n) := \mathrm{Fin}\,n\to\{E,A,B,C\},
	\qquad
	C_n := \texttt{canonicalCoveringNumber}(n)=4^n.
	\]
	
	\subsection{Von der kanonischen zur minimalen Überdeckung}
	\label{sec:canonical-to-minimal}
	
	\paragraph{Schritt 2 (geschlossen): Separation + Testmetrik.}
	In \texttt{ShellSeparation.lean}:
	\begin{itemize}
		\item diskrete Testmetrik \texttt{shellDist} ($0/1$),
		\item \texttt{ShellSeparation}: $\forall u\neq v:\; d(u,v)\ge\varepsilon$,
		\item \texttt{canonicalCellCover}: existierende kanonische Überdeckung der Größe $C_n$.
	\end{itemize}
	Bewiesen ist damit die Kette
	\[
	\text{kanonische Zellen}
	\;\Longrightarrow\;
	C_n = 4^n.
	\]
	
	\paragraph{Schritt 3 (geschlossen): minimale \emph{separierte} Überdeckung.}
	In \texttt{ShellMinimalCover.lean} wird kein $\mathbb{R}$-Infimum genommen,
	sondern ein endliches Minimum über $\mathbb{N}$, weil
	$\texttt{ShellVertex}(n)$ endlich ist:
	\[
	\texttt{ShellCover}(n):\;
	\texttt{cells}:\;\mathrm{Finset}(\mathrm{Set}\,V_n),\qquad
	\forall v\,\exists U\in\texttt{cells}:\;v\in U.
	\]
	Die kombinatorische Separationsbedingung lautet:
	jede Zelle enthält höchstens einen Vertex
	(\texttt{SeparatedShellCover.small}).
	Dann gilt die untere Schranke
	\[
	\#V_n \;\le\; \#\texttt{cells}
	\qquad
	(\texttt{separated\_cover\_lower\_bound}),
	\]
	und mit der Singleton-Überdeckung folgt
	\[
	M_n^{\mathrm{sep}} \;=\; C_n \;=\; 4^n
	\qquad
	(\texttt{minimal\_separated\_cover\_eq\_four\_pow}).
	\]
	Präzise: $M_n^{\mathrm{sep}}=C_n$, nicht pauschal $M_n=C_n$ unter
	metrischer Ball-Überdeckung.
	
	\paragraph{Skalierungsverhältnis (über $C_n$ und $M_n^{\mathrm{sep}}$, bewiesen).}
	\[
	\frac{\log C_n}{\log r_n}
	\;=\;
	\frac{\log M_n^{\mathrm{sep}}}{\log r_n}
	\;=\;
	\frac{\log 4}{\log\phi}
	\qquad
	(\texttt{fractal\_ratio\_via\_canonical\_cover},
	\texttt{fractal\_ratio\_via\_minimal\_separated\_cover}).
	\]
	\paragraph{Abschluss Phase~B.}
	Mit $r_n=\phi^n$ und $M_n^{\mathrm{sep}}=4^n$ ist das ehrliche Ergebnis
	\[
	\boxed{
		\frac{\ln M_n^{\mathrm{sep}}}{\ln r_n}
		\;=\;
		\frac{\ln 4}{\ln\phi}
	}.
	\]
	Phase~B ist damit eingefroren: formale kombinatorische Fraktalstruktur,
	ohne Anspruch auf metrische Realisierung in $\mathbb{R}^3$.
	
	\paragraph{Phase~C (Teilstand).}
	Einbettungen für $n\le 3$ und Ikosaeder-Separation sind Lean-belegt;
	metrische Minimalüberdeckung, Box-Dimension und Grenzwert $n\to\infty$
	bleiben offen; siehe Abschnitt~\ref{sec:three-phases}.
	
	\subsection{Zeitgenössische Nachbarn: Maynard, Tao, Scholze}
	\label{sec:modern-neighbors}
	
	Die EABC-Theorie operiert nicht auf dem Niveau von Gaps zwischen Primzahlen
	(Maynard--Tao) oder perfektoider Geometrie (Scholze), sondern auf einer
	\emph{diskreten mod-$12$-Klassifikation} ungerader Normkugeln.
	Der strukturelle Anschluss ist arithmetisch, nicht analytisch:
	Primzahlen $p>3$ mit $p\bmod 12\in\{1,5,7,11\}$ erzeugen privilegierte
	Defekt-Ereignisse mit quantisierten Gewichten $w_p\in\{1,5,7,11\}$.
	Die Verifikation betrifft damit Symmetrie-Fixpunkte und projektive Retraktion,
	nicht die Dichte oder Verteilung der Primzahlen selbst.
	Im historischen Bogen (Abschnitt~\ref{sec:introduction}) steht diese Arbeit
	näher bei Euler und 't~Hooft (diskrete Symmetrie, projektive Stabilisierung)
	als bei den genannten Primzahlprogrammen der Gegenwart.
	
	\paragraph{Abgrenzung: Polymath8, Maynard, Riemannsche Hypothese.}
	Das Polymath8-Projekt~\cite{zhang-polymath8} und Maynards kleine Primzahllücken~\cite{maynard-gaps}
	studieren \emph{globale} Gap-Statistiken und asymptotische Dichten der Primzahlen.
	EABC operiert auf einer \emph{lokalen} Ebene: ein einzelnes EABC-Prim $p>3$ mit
	$p\bmod 12\in\{1,5,7,11\}$ erzeugt ein quantisiertes Defekt-Label $w_p$.
	Es gibt keine Lean-Formalisation von Primzahlzwilling-Häufigkeiten,
	keine Analogie zu Zhangs beschränkten Lücken und \textbf{keinen} Bezug zur
	Riemannschen Hypothese im Formal Core.
	Die Riemann-Erwähnung in der Einleitung (Abschnitt~\ref{sec:introduction}) betrifft
	spektral-geometrische Lesarten von Isotropie, nicht analytische $\zeta$-Funktionen.
	RH bleibt explizit außerhalb des Verifikationsstands (\texttt{v1.0-formal-core}).
	
	%=============================================================================
	\section{Dualität}
	%=============================================================================
	
	Ikosaeder--Dodekaeder-Dualität: $I_{12}\leftrightarrow D_{20}$ (12 Ecken $\leftrightarrow$ 20 Flächen).
	Unter Dualitätsabbildung und Label-Pushforward gilt:
	
	\begin{lemma}[Dualitätsstabilität]
		$\Delta(M)$ am geometrischen Fixpunkt ist unter Ikosaeder$\leftrightarrow$Dodekaeder-Dualität
		invariant. Dagegen sind der Dipol $D=\sum_i w_{\sigma(i)} v_i$ und die Holonomie $H$
		\emph{nicht} dualitätsinvariant.
	\end{lemma}
	
	\smallskip\noindent
	\textit{Numerisch verifiziert in \texttt{eabc\_icosahedron\_test.py},
		\texttt{run\_core\_lemma\_tests()} (4/4 PASS, Toleranz $10^{-6}$).}
	
	%=============================================================================
	\section{Physikalische Lesart: Stella Octangula und QCD (Phase~C)}
	\label{sec:phase-c-qcd}
	%=============================================================================
	
	\textbf{Epistemischer Status.}
	Dieser Abschnitt ist \emph{kein} Teil der Lean-Verifikation (v1.0 Formal Core).
	Er skizziert eine physikalische Lesart als Brücke von der bewiesenen Tensor-Restauration
	(Theorem~\ref{thm:tensor-restore}) zur starken Wechselwirkung.
	Eine vollständige $\mathbb{R}^3$-Einbettung für alle $n$ und jede QCD-Aussage
	bleiben offen; für $n\le 3$ und Ikosaeder-Separation siehe Phase~C-Teilstand
	(Abschnitt~\ref{sec:what-we-prove}).
	
	\subsection{Das duale Tetraeder-Szenario}
	\label{sec:stella-octangula}
	
	Zwei reguläre Tetraeder, um $60^\circ$ zueinander verschoben und ineinander gestellt,
	bilden die \emph{Stella Octangula} (Kepler-Stern) --- ein klassisches Sternpolyeder
	im $\mathbb{R}^3$.
	In der hier vertretenen \emph{Lesart} (noch nicht formalisiert) wird diese Konfiguration
	als geometrisches Vorspiel für zwei komplementäre Kollektoren gelesen:
	
	\begin{enumerate}
		\item \textbf{Fraktaler Punktraum $\mathcal{T}_\infty$:}
		Ein Tetraeder-Kollektor sammelt die diskreten EABC-Primdefekte ($w_p$) über die
		Peano-Schachtelung (Abschnitt~\ref{sec:formal-verification}).
		Anschaulich: unendliche informationelle Oberfläche, asymptotisches Feldrauschen.
		\item \textbf{Messbarer Volumenraum $\mathcal{T}_V$:}
		Das um $60^\circ$ gedrehte Partner-Tetraeder fängt die Punktstruktur in einem
		endlichen Volumen ein --- analog zu einer Resonanz- oder Konfinement-Geometrie
		(Proton/Neutron als \emph{Bild}, nicht als Lean-Theorem).
	\end{enumerate}
	
	\begin{remark}[Geometrische Einbettung --- Teilstand]
		Die Ikosaeder-Repräsentanten-Einbettung auf Schale~$1$ mit Separation
		$\varepsilon=1$ ist in \texttt{IcosahedronEmbedding.lean} bewiesen
		(\texttt{icosahedronShellEmbedding\_one\_sep}).
		Die Aussage, dass die Stella Octangula in einer sphärischen Hülle eingebettet
		die zwölf Ikosaeder-Ecken \emph{global} induziert, bleibt eine motivierende
		$\mathbb{R}^3$-Lesart und ist nicht Teil von \texttt{geometric\_fixpoint\_isotropy}.
	\end{remark}
	
	\subsection{Bezug zur verifizierten Satzkette}
	
	Was \textbf{Lean beweist} (Phase~A):
	\[
	M_{\mathrm{geom}}(I_{12},\sigma_{\mathrm{geo}})=24I_3,
	\qquad
	M_{\mathrm{eff}}(R^*(K^+))=24I_3.
	\]
	Was \textbf{kombinatorisch bewiesen} ist (Phase~B, Nebenzweig):
	$\ln M_n^{\mathrm{sep}}/\ln r_n=\ln4/\ln\phi$ und $D_{EABC}=\ln4/\ln\phi<3$
	(\texttt{fractal\_dimension\_lt\_three}) --- ein Skalierungsinvarianzsatz,
	\emph{keine} bewiesene Box-Dimension und \emph{kein} QCD-Confinement-Theorem.
	
	\begin{hypothesis}[QCD-Fixpunkt-Lesart]
		\label{hyp:qcd-gauge-fix}
		In der physikalischen Lesart wirkt $R^*_{\mathrm{EABC}}$ als algebraischer
		\emph{Eich-Fixierer}: das fraktale Defektfeld $\mathcal{T}_\infty$ (Primdefekte $w_p$)
		wird durch die duale $60^\circ$-Konfiguration $\mathcal{T}_V$ in den isotropen
		Fixpunkt-Tensor projiziert:
		\[
		\mathcal{T}_\infty \;\otimes_{60^\circ}\; \mathcal{T}_V
		\;\Longrightarrow\;
		M_{\mathrm{geom}} = 24I_3.
		\]
		Die $60^\circ$-Verschiebung wird \emph{interpretiert} als Anknüpfung an die diskrete
		Phasenstruktur der Farbsymmetrie ($SU(3)$-Lesart); diese Anknüpfung ist
		\textbf{nicht} in Lean formalisiert.
	\end{hypothesis}
	
	\begin{remark}[Was diese Hypothese \emph{nicht} behauptet]
		Hypothese~\ref{hyp:qcd-gauge-fix} ist keine Lösung des QCD-Confinement-Problems
		im Sinne der Lattice-QCD oder der Yang--Mills-Millennium-Aussage.
		Sie behauptet nicht, dass $D_{EABC}<3$ die Konvergenz hadronischer Observablen
		beweist, und ersetzt keine Feynman-Diagramme oder numerische Gitterrechnung.
		Sie ordnet die \emph{bereits verifizierte} Projektionssemantik von $R^*$
		geometrisch ein --- als Forschungsprogramm Phase~C.
	\end{remark}
	
	\paragraph{Lesart der beiden Tetraeder (asymptotisch vs.\ infrarot).}
	In der QCD-Terminologie (rein interpretativ):
	\begin{itemize}
		\item $\mathcal{T}_\infty$: kurze Abstände, asymptotische Freiheit, unendliche
		virtuelle Komplexität (Peano-Stapel, Primdefekte).
		\item $\mathcal{T}_V$: große Abstände, Konfinement-\emph{Bild}, endliches Hadronenvolumen;
		$R^*$ entfernt den Rang-1-Defekt unabhängig von $v$
		(Theorem~\ref{thm:tensor-restore}, $v$-Unabhängigkeit).
	\end{itemize}
	
	\paragraph{Nächster formaler Schritt (Phase~C).}
	Um diese Lesart über Hypothese~\ref{hyp:qcd-gauge-fix} hinaus zu prüfen, braucht man:
	\begin{enumerate}
		\item explizite $\iota_n:\texttt{ShellVertex}(n)\to\mathbb{R}^3$ mit Stella-Octangula-/Ikosaeder-Geometrie,
		\item metrische Separation $\| \iota_n(u)-\iota_n(v)\|\ge\varepsilon_n$,
		\item Beweis oder Widerlegung der Verbindung zu Ball-Überdeckungen $M_n(\varepsilon_n)$.
	\end{enumerate}
	
	%=============================================================================
	\section{Superkritische Normschale und Kernladungsgrenze (Phase~C)}
	\label{sec:phase-c-supercritical}
	%=============================================================================
	
	\textbf{Epistemischer Status.}
	Dieser Abschnitt ist \emph{kein} Teil der Lean-Verifikation (v1.0 Formal Core)
	und \emph{keine} Konsequenz von Theorem~\ref{thm:tensor-restore}.
	Er formuliert eine physikalische Forschungsfrage, die später getestet oder
	verworfen werden kann --- mit derselben strengen Kennzeichnung wie
	Abschnitt~\ref{sec:phase-c-qcd}: Motivation, Beobachtung, Hypothese, offene Tests.
	
	\subsection{Etablierte Physik}
	\label{sec:supercritical-physics}
	
	Zwei Grenzen der relativistischen Quantenmechanik und QED sind hier zu trennen:
	
	\paragraph{Dirac-Grenze (idealisiert, punktförmige Kernladung).}
	Für eine punktförmige Kernladung $Z$ wird die 1s-Elektronenenergie bei
	\[
	Z\alpha \approx 1,
	\qquad
	\alpha \approx \tfrac{1}{137},
	\]
	problematisch; grob $Z\approx 137$.
	Dies ist die bekannte Dirac-Grenze: ab dort wird die relativistische
	Elektronenbeschreibung für punktförmige Kerne singular.
	
	\paragraph{Superkritische Grenze (reale Kerne, endlicher Radius).}
	Für Kerne mit endlichem Radius verschiebt sich die Schwelle auf
	\[
	Z_{\mathrm{krit}} \approx 170\text{--}173.
	\]
	Ab dort spricht man von einem superkritischen Coulomb-Feld: die innerste
	Elektronenschale (insbesondere 1s) sinkt so tief in das negative Dirac-Meer,
	dass spontane Elektron--Positron-Paarbildung möglich wird.
	In der hier verwendeten Bildsprache: die Elektronenhülle
	\emph{kollabiert} bzw.\ ist nicht mehr als normale Hülle stabil beschreibbar.
	
	\begin{remark}[Keine Primzahlarithmetik]
		Diese Aussagen stammen aus der etablierten relativistischen QED;
		sie folgen \emph{nicht} aus der EABC-Primzahlkette oder aus
		Theorem~\ref{thm:tensor-restore}.
	\end{remark}
	
	\subsection{EABC-Beobachtung}
	\label{sec:supercritical-observation}
	
	Reine arithmetische Beobachtung im EABC-Modell (kein physikalischer Beweis):
	\[
	137 \equiv 5 \pmod{12},
	\qquad
	173 \equiv 5 \pmod{12},
	\]
	also beide Primzahlen in der EABC-Klasse~$A$.
	Nach Definition von \texttt{classOf} und \texttt{primeWeight} folgt im Modell
	\[
	w_{137} = w_{173} = 5.
	\]
	\textbf{Mehr folgt zunächst nicht:}
	weder dass $137$ die Feinstrukturkonstante erklärt,
	noch dass $173$ eine Hülleninstabilität \emph{bewirbt},
	noch dass $R^*$ Vakuumpolarisation modelliert.
	
	\begin{remark}[137 vs.\ 173 im EABC-Kontext]
		$137$ liegt nahe an $\alpha^{-1}$ und markiert die idealisierte Dirac-Grenze;
		$173$ liegt nahe an $Z_{\mathrm{krit}}$ für endliche Kerne.
		Beide sind EABC-zulässige Primzahlen mit gleichem Label-Gewicht $w_p=5$ ---
		eine \emph{arithmetische} Koinzidenz, keine abgeleitete physikalische Identität.
	\end{remark}
	
	\subsection{Hypothese und offene Tests}
	\label{sec:supercritical-hypothesis}
	
	\begin{hypothesis}[Superkritische Normschalen-Schwelle]
		\label{hyp:shell-collapse}
		\boxed{\text{Hypothese: Superkritische Normschalen-Schwelle}}
		\smallskip
		
		In einer \emph{möglichen} physikalischen Lesart der EABC-Normschalen-Hierarchie
		könnte die erste superkritische Kernladungszahl $Z_{\mathrm{krit}}\approx 173$
		einer ausgezeichneten EABC-Primnormschale $p_{\mathrm{krit}}=173$ entsprechen.
		Diese Interpretation ist \textbf{nicht} aus der formalen Theorie abgeleitet
		und derzeit spekulativ.
	\end{hypothesis}
	
	\begin{remark}[Was diese Hypothese \emph{nicht} behauptet]
		\label{rem:shell-collapse-not}
		\begin{itemize}
			\item \textbf{Nicht:} $p_{\mathrm{krit}}=173$ als \emph{Resultat} des Formal Core.
			\item \textbf{Nicht:} $w_p=5$ erklärt die QED-Grenze oder $\alpha^{-1}$.
			\item \textbf{Nicht:} $R^*_{\mathrm{EABC}}$ modelliert Vakuumpolarisation
			oder Paarerzeugung aus dem Dirac-Meer.
			\item \textbf{Nicht:} $173$ folgt aus der EABC-Arithmetik oder aus
			\texttt{prime\_norm\_full\_restoration}.
		\end{itemize}
	\end{remark}
	
	\paragraph{Offene Forschungsfrage (mathematisch prüfbar).}
	Die zentrale prüfbare Frage lautet:
	\emph{Kann eine metrische Instabilität im \texttt{ShellEmbedding}-Modell auftreten,
		deren erste kritische Normschale bei einer Primzahl nahe $173$ liegt?}
	Konkret wären Kandidaten für eine solche Instabilität:
	\begin{enumerate}
		\item $\varepsilon_n\to 0$ entlang einer Peano-Schalenfolge trotz strikter Separation
		in der kombinatorischen Phase~B,
		\item Divergenz oder Kollaps der metrischen Überdeckungszahl $M_n(\varepsilon_n)$,
		\item Verletzung von \texttt{ShellEmbedding.separated} unter Skalierung
		$\lambda_{\mathrm{sphere}}$ (Abschnitt~\ref{sec:formal-verification}).
	\end{enumerate}
	Eine positive Antwort würde die Hypothese \emph{stützen}; eine negative würde sie
	\emph{widerlegen} --- ohne den Formal Core zu berühren.
	
	\paragraph{Bezug zum Lean-Code (Phase~C).}
	\texttt{ShellEmbedding.lean}: metrische Einbettung
	(\texttt{cardinalShellEmbedding\_one/two/three}, \texttt{rotateZ60};
	Ikosaeder-Separation in \texttt{IcosahedronEmbedding.lean} bewiesen).
	\texttt{DiracMonopole.lean}: vier Mathlib-U(1)-Anker, drei offene Monopol-Hypothesen.
	\texttt{ShellSeparationScale.lean}: \texttt{phiScaleAchievable\_one/two/three}
	($\varepsilon_n=\varphi^{-n}$ auf Gitter-Einbettungen).
	\texttt{ShellInstability.lean}: abstrakte Instabilitäts-API ---
	\[
	\texttt{MetricSeparationLossExists}
	\;\Rightarrow\;
	\texttt{IsFirstSeparationLoss}\,n
	\;\Rightarrow\;
	\texttt{shellPrimeCandidate}\,n = p
	\;\Rightarrow\;
	\text{Vergleich mit externem }p.
	\]
	Sanity: \texttt{not\_shellSeparationLoss\_one}
	($\neg\,\mathrm{ShellSeparationLoss}(1)$; die Definition ist nicht trivial).
	
	\begin{remark}[Keine Zahl $173$ im Formalismus]
		Die formale Phase-C-API enthält die Zahl $173$ \emph{nicht}.
		Der QED-Vergleichswert $Z_{\mathrm{krit}}\approx 173$ tritt erst als
		\emph{externer} Vergleich auf, nachdem intern eine erste Separationsverlust-Stufe
		\texttt{IsFirstSeparationLoss}\,$n$ etabliert wurde.
		\smallskip
		
		\noindent\textit{The formal Phase-C API does not contain the number $173$.
			The QED value enters only as an external comparison after an internally defined
			first separation-loss stage has been established.}
	\end{remark}
	
	\paragraph{Epistemische Einordnung (Phase~A / B / C).}
	\begin{center}
		\small
		\begin{tabular}{ll}
			\hline
			\textbf{Phase} & \textbf{Status} \\
			\hline
			A --- Formal Core & $\checkmark$ 24-Isotropy-Restoration bewiesen \\
			B --- Kombinatorik & $\checkmark$ $M_n^{\mathrm{sep}}=4^n$, Skalierungsinvarianz \\
			C --- Metrik / Physik &
			$\checkmark$ $\iota_n$ für $n\le 3$, Ikosaeder-Separation, Dirac-Anker; \\
			& $\circ$ ShellCollapse-Hypothese (Hypothese~\ref{hyp:shell-collapse}); \\
			& $\circ$ Supercritical-Threshold-Lesart; \\
			& $\circ$ QCD-/Stella-Octangula-Lesart (Hypothese~\ref{hyp:qcd-gauge-fix}); \\
			& $\circ$ Externe Analogien FSP/Kosmologie (Abschnitt~\ref{sec:phase-c-bridges}); \\
			& $\circ$ Offene Eichkorrespondenzen (Abschnitt~\ref{sec:open-problems}). \\
			\hline
		\end{tabular}
	\end{center}
	
	%=============================================================================
	\section{Physikalische Analogien: Festkörper und Kosmologie (Phase~C)}
	\label{sec:phase-c-bridges}
	%=============================================================================
	
	\textbf{Epistemischer Status.}
	Dieser Abschnitt ist \emph{kein} Teil der Lean-Verifikation (v1.0 Formal Core)
	und \emph{keine} Konsequenz von Theorem~\ref{thm:tensor-restore}.
	Er dokumentiert \textbf{externe} Anknüpfungspunkte aus der aktuellen
	Festkörperphysik und Kosmologie, die als \emph{methodische Brücken} für das
	erweiterte Bamberg-/Duo-Tetraeder-Programm dienen.
	Die numerischen Invarianten des Bamberg-Modells
	($E_0\approx -4{,}085$, $\Delta E=4$, $T_c\approx 1{,}7$)
	stehen hier \emph{parallel} zu den Lean-Invarianten
	($M_{\mathrm{geom}}=24I_3$, $M_n^{\mathrm{sep}}=4^n$, $\varepsilon_n=\varphi^{-n}$)
	und werden nicht als deren Ableitung behauptet.
	
	\subsection{Festkörperphysikalischer Anknüpfungspunkt: topologische Monopolladung}
	\label{sec:bridge-condensed-matter}
	
	Zhao und Chen~\cite{zhao-chen-2026} zeigen für Dipol-Oktupol-Doublets auf
	Pyrochlor-Gittern, dass die Wahl der Symmetrie-Untergruppe
	($U(1)$-Spin-Flüssigkeit) darüber entscheid, ob elementare Defekte (Monopole)
	eine physikalisch messbare Ladung tragen.
	In der oktupolar-dominanten Phase ist die Ladung starr null ($Q_m=0$),
	in der dipolar-dominanten Phase friert ein endlicher Ganzzahl-Quantenwert ein
	($Q_m\neq 0$).
	
	\begin{remark}[Algebraische Lesart (Bamberg-Programm, nicht formalisiert)]
		Im Bamberg-Modell wird dieser Übergang \emph{interpretiert} als geometrische
		Frustration beim Wechsel von der quaternionischen Unterstruktur ($\mathbb{H}$)
		zum oktionischen Gitter ($\mathbb{O}$).
		Der nackte Austauschterm $J/4$ aus dem Hantel-Modell
		\[
		v_0 = \frac{J}{4} + \frac{1}{12}\left(4\sqrt{\tfrac{2}{3}}-1\right)
		\frac{\mu_0 \tilde{g}^2 \mu_B^2}{4\pi r_{nn}^3}
		\]
		wird formal mit dem ungestörten Clifford-Kerndruck $P_0/4=16/4=4{,}0$
		in Parallelität gesetzt; das kontinuierliche Korrekturglied höherer Ordnung
		wird als physikalisches Äquivalent zum KZM-Fehlerfaktor
		$\delta\approx +0{,}145954$ gelesen.
		\textbf{Diese Isomorphie ist eine Arbeitshypothese, kein Lean-Theorem.}
	\end{remark}
	
	\begin{hypothesis}[Hurwitz-Orbit und emergente Eichladung]
		\label{hyp:hurwitz-charge}
		Wir postulieren, dass die Wahl des Hurwitz-Orbits in der algebraischen Struktur
		die quantisierte Natur der emergenten Eichladung bestimmt.
		Analog zu den Befunden von Zhao et al.~\cite{zhao-chen-2026}, bei denen die
		topologische Ladung $Q_m$ starr an die dipolare Symmetrieachse gekoppelt ist,
		generiert die Nichtassoziativität des $E_8$-Vakuums (in der Bamberg-Lesart)
		eine weitreichende Coulomb-Wechselwirkung, welche ein rein lokales Screening
		blockiert.
	\end{hypothesis}
	
	\begin{remark}[Was Hypothese~\ref{hyp:hurwitz-charge} \emph{nicht} behauptet]
		Sie behauptet nicht, dass $Q_m\sim 10^{-5}q_D$ aus Dipol-Spineisen
		\emph{bewiesen} aus $M_{\mathrm{geom}}=24I_3$ folgt,
		noch dass $R^*_{\mathrm{EABC}}$ ein Gitter-Eichfeld auf Pyrochloren modelliert.
		Sie formuliert eine prüfbare algebraisch-physikalische Korrespondenz als
		Forschungsprogramm.
	\end{remark}
	
	\subsection{Kosmologischer Anknüpfungspunkt: relativistische Gitterinstabilität}
	\label{sec:bridge-cosmology}
	
	Alexander, Temple und Vogler~\cite{alexander-temple-vogler-2026} zeigen über
	eine Reihenentwicklung der Einstein-Euler-Gleichungen in selbstähnlichen
	Schwarzschild-Variablen (STV), dass das flache, kritische Friedmann-Universum
	($k=0$) ein instabiler Sattelpunkt ($SM$) am Urknall ist:
	jede generische, unterdichte Fluktuation muss von diesem Zustand wegevolvieren.
	In den autonomen STV-ODEs tritt bei Ordnung $n=2$ (einem $4\times4$-System)
	ein singulärer negativer Eigenwert $\lambda_{B2}=-1/3$ auf;
	alle höheren Eigenwerte ab $n\ge 3$ sind streng positiv.
	
	\begin{remark}[Numerische Resonanz im Bamberg-Programm]
		In dieser Lesart wird die im Bamberg-Modell beobachtete Anregungsenergie
		$\Delta E=4{,}0000$ als \emph{strukturelle} Entsprechung zur kritischen
		Ordnung $n=2$ nach Alexander et al.~\cite{alexander-temple-vogler-2026}
		geführt: die Emergenz von $\lambda_{B2}<0$ bricht die makroskopische
		Selbstähnlichkeit über intermediäre Zeitskalen auf und erzeugt transiente,
		anomale Beschleunigung (phenomenologisch nahe $Q\approx 0{,}425$ in der
		Dunkle-Energie-Lesart), bevor die Traktatorien asymptotisch gegen den
		Minkowski-Ruhpunkt $M$ ($t\to\infty$) zerfallen.
		\textbf{Keine dieser Kosmologie-Aussagen ist in Lean formalisiert.}
	\end{remark}
	
	\begin{hypothesis}[Primzahlzwillinge und STV-Attraktor]
		\label{hyp:prime-twin-stv}
		Die numerische Robustheit der Invariante $\Delta E=4$ im Bamberg-Programm
		korrespondiert strukturell mit der Eigenwert-Hierarchie der Einstein-Euler-Gleichungen
		bei kritischer Ordnung $n=2$.
		Die daraus resultierende Feldstauchung wird als topologische Stringspannung
		gelesen, welche das Confinement einer unendlichen Primzahlzwillingskette im
		minimalen Phasenabstand $\Delta=2$ \emph{erzwingen könnte}.
	\end{hypothesis}
	
	\subsection{Synthese: Invarianten vs.\ Lesarten}
	
	Unabhängig davon, ob man die Gitterverzerrung festkörperphysikalisch
	(Spin-Eis) oder kosmologisch (ART) interpretiert, bleibt im Bamberg-Programm
	$\Delta E=4$ als geometrische Kennzahl des Modells bestehen.
	Im Formal Core bleibt parallel $M_n^{\mathrm{sep}}=4^n$ und
	$\varepsilon_n=\varphi^{-n}$ kombinatorisch/metrisch getrennt von jeder
	physikalischen Identifikation.
	Abschnitt~\ref{sec:what-we-prove} liefert die bewiesenen Fakten;
	Abschnitt~\ref{sec:phase-c-bridges} die externen Analogien;
	Abschnitt~\ref{sec:open-problems} die präzisen mathematischen Meilensteine.
	
	%=============================================================================
	\section{Offene mathematische Probleme und Eichkorrespondenzen}
	\label{sec:open-problems}
	%=============================================================================
	
	\textbf{Epistemischer Status.}
	Die Robustheit der invarianten Kennzahlen
	($E_0\approx -4{,}085$, $\Delta E=4$, $T_c\approx 1{,}7$ im Bamberg-Programm;
	$M_{\mathrm{geom}}=24I_3$, $M_n^{\mathrm{sep}}=4^n$ im Formal Core)
	erfordert die Formulierung fundamentaler mathematischer Anschlussfragen
	an die Community.
	Keine der folgenden Fragen ist im gegenwärtigen Lean-Release beantwortet.
	
	\begin{enumerate}
		\item \textbf{Topologische Ladungskorrespondenz.}
		Lässt sich die aus dem Hantel-Modell emergente Monopolladung
		$Q_m\sim 10^{-5}q_D$ von Dipol-Spineisen~\cite{zhao-chen-2026}
		rigoros als algebraischer Index (Windungszahl) auf die diskreten
		Hurwitz-Orbits abbilden?
		\item \textbf{Stabilität des Attraktors.}
		Konvergiert die unendliche Kette der Primzahlzwillinge im Gitterabstand
		$\Delta=2$ mathematisch zwingend gegen die unterdichte Familie $\mathcal{F}$
		der Einstein-Euler-Gleichungen, um den von Alexander et al.~\cite{alexander-temple-vogler-2026}
		nachgewiesenen asymptotischen Zerfall gegen den universellen Minkowski-Zustand
		$M$ ($t\to\infty$) energetisch zu garantieren?
		\item \textbf{Nichtassoziativität als Eichfeld.}
		Kann der Übergang von der lokalen algebraischen Frustration des
		Duo-Tetraeders zu einer kontinuierlichen Berry-Phase differentialgeometrisch
		als holonome Verbindung (Connection) konstruiert werden, die das
		phänomenologische Confinement-Potential
		$H_{ab}\sim \mathrm{Tr}(\sigma^a R\,\sigma^b R^\dagger)$
		abbildet?
	\end{enumerate}
	
	\begin{remark}[Gutachter-Einordnung]
		Eine positive Antwort auf eine dieser Fragen würde das erweiterte
		Forschungsprogramm stützen; eine negative würde die jeweilige Lesart
		einschränken oder verwerfen --- \emph{ohne} den Formal Core
		(24-Isotropy-Restoration) zu berühren.
	\end{remark}
	
	\subsection{Geometrie-Blueprint: Morley, Marion--Walter, Ptolemy}
	\label{sec:geometry-scaffold}
	
	\textbf{Epistemischer Status.}
	Das Modul \texttt{GeometryScaffold.lean} überträgt die Tao/PFR-Methodik
	(Abschnitt~\ref{sec:what-we-prove}, \texttt{PhaseCBlueprint.lean}) auf
	klassische euklidische Konstruktionen im EABC-Kontext.
	Es beweist weder Morleys Dreieck noch die Marion--Walter-Flächenformel;
	alle Ziele sind als offene \texttt{Prop}s markiert, analog zu
	\texttt{DiracMonopole.lean} (Abschnitt~\ref{sec:dirac-monopole}).
	
	\paragraph{Drei Beweisstruktur-Vorlagen (Phase~C).}
	\begin{enumerate}
		\item \textbf{Morley (Winkeltrisierung):}
		\texttt{AngleTrisectionScaffold} + \texttt{MorleyEquilateralHypothesis} ---
		gleichseitiges Dreieck aus Trisektions-Scaffold (offen).
		\item \textbf{Marion--Walter (Hexagonfläche):}
		\texttt{MarionWalterHexagon} + \texttt{MarionWalterAreaHypothesis} ---
		Hexagonfläche $=\tfrac{1}{10}$ der Dreiecksfläche (offen).
		\item \textbf{Ptolemy (Viereck):}
		\texttt{CrossedPtolemyHypothesis} (Bowtie-Variante, offen) und
		\texttt{PtolemyShellQuadrupleHypothesis} (Ikosaeder-Repräsentanten
		$A$--$E$--$C$--$B$ unter \texttt{icosahedronShellEmbedding\_one}, offen).
	\end{enumerate}
	
	\begin{table}[h!]
		\centering
		\small
		\begin{tabular}{@{}lll@{}}
			\hline
			\textbf{Meilenstein} & \textbf{Lean} & \textbf{Status} \\
			\hline
			Morley gleichseitig & \texttt{morley\_equilateral} & rot \\
			Marion--Walter $1/10$ & \texttt{marion\_walter\_hexagon} & rot \\
			Gekreuztes Ptolemy (Bowtie) & \texttt{crossed\_ptolemy\_bowtie} & rot \\
			Ptolemy Ikosaeder-Viereck & \texttt{ptolemy\_shell\_quadruple} & rot \\
			Mathlib Ptolemy (kosphärisch) & \texttt{mathlib\_ptolemy} & grün \\
			\hline
		\end{tabular}
		\caption{Geometrie-Meilensteine in \texttt{PhaseCBlueprint}/\texttt{VerificationCertificate}.}
		\label{tab:geometry-milestones}
	\end{table}
	
	\begin{remark}[Empfohlener nächster Schritt]
		Der \textbf{Ptolemy-Pfad} ist der strukturell nächste grüne Schritt:
		\texttt{milestone\_green\_mathlib\_ptolemy} importiert
		\texttt{mul\_dist\_add\_mul\_dist\_eq\_mul\_dist\_of\_cospherical};
		\texttt{PtolemyShellQuadrupleHypothesis} fixiert bereits das Ikosaeder-Viereck
		(\texttt{icosaRepQuadruple}, paarweise verschiedene Repräsentanten-Klassen
		\texttt{icosaRepQuadruple\_pairwise\_ne}) und \texttt{shellDist3}-Produkte als prüfbares Ziel.
		Eine Verifikation würde die metrische Schicht (Abschnitt~\ref{sec:phase-c-bridges})
		mit dem kombinatorischen Ikosaeder-Fixpunkt $M_{\mathrm{geom}}=24I_3$ verbinden ---
		ohne Morley- oder Marion--Walter-Beweise vorwegzunehmen.
	\end{remark}
	
	\subsection{12-Kuss-Geometrie und Schalen-Durchlass}
	\label{sec:kissing-gap}
	
	\textbf{Epistemischer Status.}
	Das Modul \texttt{KissingGap.lean} formalisiert die klassische Ikosaeder-Kuss-Konfiguration
	(1 zentrale + 12 äußere Kugeln gleichen Radius~$R$) als \emph{geometrische Analogie}
	zur EABC-Schalenmetrik (Phase~C). Bewiesen sind die Konstantenidentitäten und der
	Verknüpfungsanker \texttt{icosaShortEdge\_dist}; offen bleiben
	\texttt{KissingGapHypothesis} und \texttt{ShellPassageHypothesis}.
	
	\paragraph{EABC-Zuordnung (Potential vs.\ Kraft).}
	\begin{center}
		\small
		\begin{tabular}{@{}ll@{}}
			\hline
			\textbf{12-Kuss-Objekt} & \textbf{EABC-Lesart} \\
			\hline
			Zentrale Kugel & geometrischer Fixpunkt / zugewiesene Prim-Schale \\
			12 äußere Kugeln & 12-Kuss-Richtungen, \texttt{icosaVert}, \texttt{ShellEmbedding} \\
			Kuss-Konfiguration (global) & globales Potential (\texttt{M\_geom}, 12-Kuss-System) \\
			Lücke / Durchlass (lokal) & metrische Zwangsbedingung (\texttt{ShellSeparationLoss}-Analog) \\
			13.\ Vertex durch den Ring & offene Schalen-Durchlass-Hypothese \\
			\hline
		\end{tabular}
	\end{center}
	
	\paragraph{Numerische Verhältnisse} (gleiche Radien $R_c=R_o=R$, Lean \texttt{KissingGap}):
	\begin{table}[h!]
		\centering
		\small
		\begin{tabular}{@{}lll@{}}
			\hline
			\textbf{Größe} & \textbf{Lean} & \textbf{Wert} \\
			\hline
			Kanten/Umkreis $a/\rho$ & \texttt{icosaEdgeToCircumRatio} & $4/\sqrt{10+2\sqrt5}\approx 1{,}052$ \\
			Lückenfraktion $g/R$ & \texttt{kissingGapFraction} & $2(a/\rho-1)\approx 10{,}3\,\%$ \\
			Max.\ Durchlass $r/R$ & \texttt{maxPassableBallFraction} & $a/\rho-1\approx 5{,}15\,\%$ \\
			Schwelle $R_o/R_c$ & \texttt{outerShrinkForPassage} & $\approx 0{,}357$ ($\approx 64\,\%$ Schrumpfung) \\
			\hline
		\end{tabular}
		\caption{12-Kuss-Konstanten; \texttt{equal\_radii\_no\_passage} beweist: bei $R_c=R_o$ kein Durchlass.}
		\label{tab:kissing-gap}
	\end{table}
	
	\begin{remark}[Uniforme Skalierung]
		Bei gleichen Radien ist $g/R\approx 10{,}3\,\%$ und $2R_o>g$ --- ein gleich großer
		13.\ Körper passiert den Ring nicht (\texttt{equal\_radii\_no\_passage}).
		Uniforme Skalierung ändert die Verhältnisse nicht; Durchlass erfordert
		$R_o/R_c\le\texttt{outerShrinkForPassage}$ (Hypothese \texttt{ShellPassageHypothesis}).
	\end{remark}
	
	\begin{table}[h!]
		\centering
		\small
		\begin{tabular}{@{}lll@{}}
			\hline
			\textbf{Meilenstein} & \textbf{Lean} & \textbf{Status} \\
			\hline
			12-Kuss-Lücke vs.\ Schale & \texttt{kissing\_gap} & rot \\
			13.\ Schalen-Durchlass & \texttt{shell\_passage} & rot \\
			\texttt{icosaShortEdge\_dist} (Anker) & \texttt{icosaShortEdge\_dist} & grün \\
			\hline
		\end{tabular}
		\caption{Kuss-Meilensteine in \texttt{PhaseCBlueprint}/\texttt{VerificationCertificate};
			vgl.\ Abschnitt~\ref{sec:geometry-scaffold}.}
		\label{tab:kissing-milestones}
	\end{table}
	
	\subsection{Dirac-Monopol, Spin-Eis und Lean-Anker (2024--2026)}
	\label{sec:dirac-monopole}
	
	\textbf{Epistemischer Status.}
	Die folgenden Aussagen sind \emph{Hypothesen} bzw.\ \emph{Bemerkungen} zur
	physikalischen Anschlussliteratur. Lean beweist weder $Q_m = n \cdot (h/2e)$
	noch effektive Monopolladungen in Spin-Eis-Materialien.
	
	\paragraph{Neue Literatur (2024--2026).}
	\begin{itemize}
		\item \textbf{Dirac-Schwinger-Quantisierung emergenter Monopole:}
		Farhana, Saccone \& Ward~\cite{ward-farhana-saccone-2025}
		(arXiv:2501.04704, \emph{Phys.\ Lett.\ B}) argumentieren, dass die in
		künstlichem Spin-Eis gemessenen effektiven Ladungen innerhalb der Fehler
		die Bedingung $eg/(4\pi\hbar)=n/2$ erfüllen.
		\item \textbf{Dirac-Strings in Raum-Zeit:}
		Huster-Zapke \& Holdsworth~\cite{huster-holdsworth-2025}
		(arXiv:2505.14811) definieren Dirac-Strings entlang Monopol-Trajektorien
		und zeigen deren Beitrag zum magnetischen Rauschsignal.
		\item \textbf{Dipolar vs.\ oktupolar (SET):}
		Zhao \& Chen~\cite{zhao-chen-2026} (arXiv:2505.07805) schlagen vor,
		dipolare und oktupolare Spin-Eis-Regime klassisch über endliche bzw.\
		verschwindende effektive Monopolladung $q_{\mathrm m}$ zu unterscheiden;
		Smith \emph{et al.}~\cite{smith-gaulin-2025} (PRX~15, 021033) finden
		für Ce$_2$Zr$_2$O$_7$ oberhalb $T=0{,}05$\,K keine oktupolaren Streusignale.
		\item \textbf{Experimentelle Dynamik:}
		Tang \emph{et al.}~\cite{tang-gegenwart-2025} (arXiv:2509.18422) berichten
		magnetischen Transport via Spin-Seebeck-Effekt in Dy$_2$Ti$_2$O$_7$;
		van den Berg \emph{et al.}~\cite{van-den-berg-ladak-2025}
		(arXiv:2511.04877) visualisieren divergente Monopolfelder in 3D-künstlichem
		Spin-Eis; Gao \emph{et al.}~\cite{gao-dai-2026} (arXiv:2601.03202) trennen
		spektroskopisch Photonen und Spinons in Ce$_2$Zr$_2$O$_7$.
	\end{itemize}
	
	\begin{hypothesis}[Topologische Ladungskorrespondenz (Präzisierung)]
		\label{hyp:monopole-winding}
		Die in Frage~(1) formulierte Abbildung „effektive Monopolladung $\to$ algebraischer
		Index auf Hurwitz-Orbits'' bleibt \emph{offen}. Die jüngste Spin-Eis-Literatur
		liefert experimentelle und modellhafte Anhaltspunkte für diskrete bzw.\
		regime-abhängige effektive Ladungen~\cite{zhao-chen-2026,ward-farhana-saccone-2025},
		belegt aber keine Isomorphie zu den EABC-Hurwitz-Orbits.
	\end{hypothesis}
	
	\begin{remark}[Vier Lean-4-Adaptationen (Mathlib)]
		Im Modul \texttt{DiracMonopole.lean} sind vier Mathlib-Bausteine als Meilensteine
		verdrahtet (grün = importiert, nicht EABC-spezifisch):
		\begin{enumerate}
			\item \texttt{Circle.periodic\_exp} --- $2\pi$-Periodizität der U(1)-Phase;
			\item \texttt{AddCircle.homeomorphCircle'} --- $S^1 \simeq \mathrm{AddCircle}(2\pi)$;
			\item \texttt{Path.Homotopic.equivalence} --- Homotopie-Äquivalenz auf Pfaden;
			\item \texttt{Path.refl} / \texttt{FundamentalGroup} (Scaffold \texttt{loopClass}) ---
			Schleifen- bzw.\ $\pi_1$-Vorbereitung.
		\end{enumerate}
		Drei offene Props (\texttt{MonopoleHurwitzWindingHypothesis},
		\texttt{DiracSchwingerEmergentHypothesis},
		\texttt{BerryHolonomyConnectionHypothesis}) markieren die Forschungslücke
		explizit als \emph{unbewiesen}.
	\end{remark}
	
	\begin{remark}[AdS/CFT-Bulk (kein Lean-Gegenstück)]
		Eine holographische AdS/CFT-Bulk-Formalisation für die EABC-Peano-Schalen
		ist nicht Teil des gegenwärtigen Lean~4-Formalismus.
		Entsprechende Literaturanker bleiben explizit außerhalb des Verifikationsstands;
		es existiert kein Lean-Modul für Bulk-Grenz-Korrespondenz.
	\end{remark}
	
	\subsection{Potentiale vs.\ Kräfte als epistemische Faktorisierung}
	\label{sec:potential-vs-force}
	
	Der methodische Rat „Potentiale statt Kräfte'' ist im EABC-Programm bereits als
	\textbf{zwei logisch unabhängige Stränge} angelegt (Abschnitt~\ref{sec:what-we-prove}):
	globale Zertifikate gegenüber lokaler Diagnostik.
	Die Zwei-Ebenen-Architektur $M_{\mathrm{geom}}$ vs.\ $M_{\mathrm{eff}}$
	(Abschnitt~\ref{sec:two-levels}) und die Idempotenz von $R^*_{\mathrm{EABC}}$
	(Lemma~\ref{lem:idempotence}) spiegeln dieselbe Struktur im Eichtheorie-Vokabular:
	ein korrekt gewählter Retraktionsschritt stabilisiert den Fixpunkt,
	ohne die gestörte Matrix dynamisch zu glätten.
	Lean beweist weder Maxwell-Dynamik noch Spin-Eis-Monopole; die EM-/Topo-Sprache
	ist eine dokumentierte Brücke zu Phase~C (Abschnitt~\ref{sec:formal-verification}).
	
	\begin{table}[h!]
		\centering
		\small
		\begin{tabular}{@{}p{0.18\linewidth}p{0.36\linewidth}p{0.34\linewidth}@{}}
			\hline
			\textbf{Lesart} & \textbf{Potential-artig (global)} & \textbf{Kraft-artig (lokal)} \\
			\hline
			Phase~A &
			$w_p$, $R^*_{\mathrm{EABC}}$, $M_{\mathrm{geom}}=24I_3$ &
			$\Delta(M^+)=w_p$, Rang-1-Nicht-Isotropie \\
			Phase~C (metrisch) &
			$\mathrm{HasShellEmbedding}(n)$, $\varepsilon_n>0$,
			Mathlib-Ptolemy-Anker &
			$\mathrm{ShellSeparationLoss}(n)$,
			\texttt{PtolemyShellQuadrupleHypothesis} \\
			Phase~C (topologisch) &
			U(1)-Anker, \texttt{GaugePotential}, \texttt{HolonomyDatum} &
			\texttt{EmergentMonopoleCharge} ($q_{\mathrm m}$), offene Hypothesen \\
			\hline
		\end{tabular}
		\caption{Epistemische Faktorisierung: Zertifikate vs.\ Messgrößen
			(keine Physikidentität; vgl.\ Tabelle~\ref{tab:three-epistemic-layers}).}
		\label{tab:potential-force}
	\end{table}
	
	\paragraph{Operationalisierungsleiter (Level~0/1/2).}
	Die Lean-Datei \texttt{DiracMonopole.lean} dokumentiert dieselbe Taxonomie
	explizit im Modulkopf; im Manuskript lesen wir sie als drei epistemische Stufen:
	
	\begin{enumerate}
		\item[\textbf{Level~0 --- Bewiesen / Scaffold.}]
		$w_p$ eindeutig (\texttt{primeWeight}, \texttt{classOf});
		$R^*_{\mathrm{EABC}}$ konstruktiv (\texttt{rStar}\,$=$\,\texttt{rEabc});
		$M_{\mathrm{geom}}=24I_3$ (\texttt{geometric\_fixpoint\_isotropy});
		$\Delta(M_{\mathrm{eff}}(K^+))=w_p$ bei $\|v\|=1$
		(\texttt{delta\_prime\_defect\_at\_fixpoint},
		\texttt{rankOne\_defect\_not\_isotropic}) --- \emph{Parallelpfad}, keine
		Voraussetzung von \texttt{prime\_norm\_full\_restoration};
		$\neg\,\mathrm{ShellSeparationLoss}(n)$ für $n\le 3$
		(\texttt{milestone\_green\_no\_loss\_one/two/three},
		via \texttt{cardinalShellEmbedding\_one/two/three});
		vier grüne U(1)-Meilensteine
		(\texttt{milestone\_green\_u1\_period},
		\texttt{milestone\_green\_circle\_homeomorph},
		\texttt{milestone\_green\_loop\_homotopy},
		\texttt{milestone\_green\_fundamental\_group};
		Abschnitt~\ref{sec:dirac-monopole});
		grüner Mathlib-Ptolemy-Anker
		(\texttt{milestone\_green\_mathlib\_ptolemy};
		Abschnitt~\ref{sec:geometry-scaffold}).
		
		\item[\textbf{Level~1 --- API ohne \texttt{sorry}.}]
		Reine Typisierung in \texttt{DiracMonopole.lean}, ohne Gleichheitsbehauptungen
		an die Physik:
		\texttt{GaugePotential} (U(1)-Phase und optionale Windungszahl),
		\texttt{HolonomyDatum} (Pfad plus Endphase auf der $\pi_1$-Vorbereitung),
		\texttt{EmergentChargeCertificate} (Paarung von $w_p$ mit
		\texttt{DiracSchwingerClass}),
		\texttt{BerryPhasePlaceholder} (projiziert \texttt{HolonomyDatum} auf die
		gespeicherte Endphase).
		Diese Strukturen sind \emph{Platzhalter-API}: sie fixieren Datentypen für eine
		spätere Korrespondenz „algebraischer Index $\approx$ Windungszahl
		$\approx$ effektive Monopolladung'', beweisen aber keine Identität.
		
		\item[\textbf{Level~2 --- Externe Hypothesen.}]
		Offene \texttt{Prop}s:
		\texttt{MonopoleHurwitzWindingHypothesis},
		\texttt{DiracSchwingerEmergentHypothesis},
		\texttt{BerryHolonomyConnectionHypothesis}
		(sowie \texttt{DipoleOctupoleMonopoleChargeHypothesis} für das
		Zhao--Chen-Regime $q_{\mathrm m}\neq 0$ vs.\ $q_{\mathrm m}=0$;
		Abschnitt~\ref{sec:dirac-monopole}, Hypothese~\ref{hyp:monopole-winding}).
	\end{enumerate}
	
	\paragraph{Lesart der Level-1-Typen.}
	\texttt{GaugePotential} und \texttt{HolonomyDatum} modellieren
	\emph{integrative} Daten: eine globale U(1)-Phase bzw.\ eine Schleifenholonomie
	als Vorbereitung für Windungszahl-Argumente.
	\texttt{EmergentChargeCertificate} koppelt das arithmetische Primlabel $w_p$
	an eine halbganzzahlige Dirac--Schwinger-Klasse --- strukturell, nicht als
	bewiesene Quantisierung $Q_m=n(h/2e)$.
	\texttt{BerryPhasePlaceholder} ist bewusst minimal: es extrahiert nur die
	Endphase aus \texttt{HolonomyDatum}; eine differentialgeometrische
	Verbindungsgleichung bleibt Hypothese
	(\texttt{BerryHolonomyConnectionHypothesis}, offenes Problem~(3) oben).
	
	\begin{remark}[Hauptpfad vs.\ Parallelpfad]
		Die Restauration $M_{\mathrm{eff}}(R^*(K^+))=24I_3$ hängt \emph{nicht} von
		$\Delta$ oder $v$ ab; $\Delta=w_p$ ist Mess-Diagnostik auf dem Parallelpfad
		(Abschnitt~\ref{sec:what-we-prove}).
		Das ist die präzise EABC-Lesart von „Potentiale tragen die Wahrheit,
		Kräfte sind Beobachtung'':
		$M_{\mathrm{geom}}$ und $R^*$ sind die beweisbare globale Semantik;
		$\Delta$, $\mathrm{ShellSeparationLoss}$ und $q_{\mathrm m}$ sind lokale
		oder regime-abhängige Diagnostik --- ohne logische Rückwirkung auf den
		Formal Core.
	\end{remark}
	
	\begin{remark}[Epistemische Trennung]
		\textbf{Bewiesen} sind alle Level-0-Aussagen der Operationalisierungsleiter.
		\textbf{Scaffold} sind die Level-1-Typen: maschinenlesbar, \texttt{sorry}$=0$,
		aber ohne physikalische Gleichungen.
		\textbf{Hypothese} sind alle Level-2-\texttt{Prop}s und die Spin-Eis-Literatur
		(Abschnitt~\ref{sec:dirac-monopole}).
		Eine Verwechslung dieser Stufen würde den Verifikationsstand des Releases
		\texttt{v1.0-formal-core} verfälschen.
	\end{remark}
	
	\subsection{Flyspeck-Zertifikate und LP-Trennung (Hales)}
	\label{sec:flyspeck-certificates}
	
	Thomas Hales' Flyspeck-Projekt~\cite{flyspeck} formalisiert den
	Kepler-Beweis in drei \emph{getrennten} Schichten:
	\begin{enumerate}
		\item \textbf{Lean/Isabelle/HOL Light} --- deduktive Beweistexte;
		\item \textbf{Python/C++} --- numerische Intervallarithmetik und Ungleichheitsprüfungen;
		\item \textbf{externes LP} --- lineare Programme für Packungs-Ungleichungen,
		unabhängig vom Beweistext verifizierbar.
	\end{enumerate}
	
	\paragraph{Was EABC übernimmt.}
	\begin{itemize}
		\item \textbf{Dreischicht-Trennung:} Lean beweist Struktur (Phase~A/B/C);
		\texttt{eabc\_icosahedron\_test.py} liefert numerische Zertifikate
		(Toleranz $10^{-6}$); \texttt{VerificationCertificate.lean} dokumentiert
		Metadaten (\texttt{coreLemmaCertificates}, \texttt{phaseCMilestones}).
		\item \textbf{Grün/rot-Disziplin:} \texttt{PhaseCBlueprint.milestoneProp} markiert
		bewiesene vs.\ offene Ziele analog zum PFR-Blueprint~\cite{tao-pfr}.
		\item \textbf{Kein Ersetzen:} \texttt{all\_core\_certificates\_verified} belegt nur
		den \emph{Status} der Python-Tests; es behauptet keine logische Äquivalenz
		zu Lean-Theoremen.
	\end{itemize}
	
	\paragraph{Was EABC nicht portiert.}
	\begin{itemize}
		\item \textbf{Kein externes LP:} metrische Untergrenzen (Ball-Überdeckungen,
		Separationsradien für $n>3$) sind nicht als Flyspeck-artige LP-Zertifikate
		formalisiert; die LP/Farkas-Schicht bleibt prospektiv
		(vgl.\ Abschnitt~\ref{sec:open-problems}).
		\item \textbf{Keine tausende Ungleichungen:} der Formal Core hat wenige
		algebraische Hauptsätze; numerische Zertifikate betreffen nur vier Kern-Lemmata.
		\item \textbf{Keine Packungsgeometrie:} Flyspeck's Kugelpackungs-LPs haben kein
		direktes EABC-Gegenstück; die Analogie ist \emph{methodisch}, nicht inhaltlich.
	\end{itemize}
	
	\begin{remark}[Brücke zu Phase-C-Brücken]
		Die Flyspeck-Methodik trennt Beweis von Rechnung; die physikalischen Brücken
		(Abschnitt~\ref{sec:phase-c-bridges}) und Monopol-Hypothesen
		(Abschnitt~\ref{sec:dirac-monopole}) bleiben bewusst außerhalb aller
		Zertifikats-Schichten --- sie sind Lesarten, keine Lean-Theoreme.
	\end{remark}
	
	%=============================================================================
	\section{Renormierungskette}
	%=============================================================================
	
	Die diskrete Kette (ohne Primschritt) lautet
	\[
	K_n \;\longrightarrow\; K_{n+1}
	\;=\;
	\underbrace{\text{Volumen-Renormierung}}_{\lambda_{\mathrm{sphere}}}
	\;\circ\;
	\underbrace{R^*_{\mathrm{EABC}}}_{P_I \circ R_{\mathrm{EABC}}}
	\;\circ\;
	\underbrace{D_{20}}_{\text{Dual}}
	\;\circ\;
	\underbrace{I_{12}}_{\text{Ikosaeder}}
	\;\circ\;
	\underbrace{\text{12-Kuss}}_{\text{Einheitsrichtungen}}.
	\]
	
	\paragraph{Erweiterung mit Prim-Erzeugungsschritt (Abschnitt~\ref{sec:prime-step}).}
	Wenn $P(2n)=p$ prim, setze
	\[
	K_n^+ \;:=\; \mathrm{Extend}(p)\,K_n,
	\qquad
	K_{n+1}
	\;=\;
	\lambda_{\mathrm{sphere}}
	\circ R^*_{\mathrm{EABC}}
	\circ \mathrm{Extend}(P)
	\;(K_n),
	\]
	wobei $\mathrm{Extend}(p)$ die Konfiguration um die Prim-Normkugel mit
	Label $\mathrm{EABC}(p)$ erweitert.
	Äquivalent:
	\[
	K_n \;\longrightarrow\; K_n^+ \;\longrightarrow\;
	R^*_{\mathrm{EABC}}(K_n^+) \;\longrightarrow\; K_{n+1}.
	\]
	
	\begin{definition}[Volumen-Renormierung]
		Für äußeren Radius $R>0$ setze $R_{\mathrm{out}}=R(1+2\pi)$ und
		\[
		\lambda_{\mathrm{sphere}}
		\;=\;
		\left(\frac{V_0}{V_{\mathrm{out}}}\right)^{\!1/3}
		\;=\;
		\frac{1}{1+2\pi},
		\]
		mit $V_0=\tfrac{4}{3}\pi R^3$, $V_{\mathrm{out}}=\tfrac{4}{3}\pi R_{\mathrm{out}}^3$.
		Die Vertex-Skalierung erfolgt isotrop: $v_i\mapsto \lambda_{\mathrm{sphere}}\,v_i$.
	\end{definition}
	
	\begin{lemma}[Volumen-Kompatibilität]
		Isotrope Skalierung mit $\lambda_{\mathrm{sphere}} = 1/(1+2\pi)$ multipliziert alle
		Eigenwerte von $M$ mit $\lambda_{\mathrm{sphere}}^2$; $\Delta(M)$ bleibt erhalten.
	\end{lemma}
	
	\smallskip\noindent
	\textit{Numerisch verifiziert in \texttt{eabc\_icosahedron\_test.py},
		\texttt{run\_core\_lemma\_tests()} (4/4 PASS, Toleranz $10^{-6}$).}
	
	%=============================================================================
	\section{Numerische Befunde}
	%=============================================================================
	
	Implementierung in \texttt{eabc\_icosahedron\_test.py}:
	
	\begin{itemize}
		\item $\varepsilon=0$: Nach $R_{\mathrm{EABC}}$ bzw.\ $R^*_{\mathrm{EABC}}$
		gilt $\Delta(M)\approx 0$ (Fixpunkt erreicht).
		\item $\varepsilon>0$: Gestörte Geometrie ändert mod-$12$-Klassifikation;
		$R_{\mathrm{EABC}}$ allein reicht nicht --- gekoppelter Operator $R^*_{\mathrm{EABC}}$
		(Pullback + Label-Projektion) konvergiert zuverlässiger zu $\Delta(M)\approx 0$.
		\item Gekoppelt vs.\ $R_{\mathrm{EABC}}$ allein: Bei $\varepsilon>0$ liefert
		$R^*_{\mathrm{EABC}}$ niedrigere mittlere $\Delta(M)$ und höhere Fixpunkt-Konvergenzrate.
		\item Volumen-Schritt: $\lambda_{\mathrm{sphere}}=1/(1+2\pi)$; $\Delta(M)$ bleibt
		nach isotroper Skalierung null, wenn es vorher null war.
	\end{itemize}
	
	\subsection{Numerische Verifikation}
	
	Die vier Kern-Lemmata werden von \texttt{run\_core\_lemma\_tests()} in
	\texttt{eabc\_icosahedron\_test.py} geprüft (Toleranz $10^{-6}$):
	
	\begin{center}
		\begin{tabular}{clc}
			\hline
			Lemma & Testschlüssel & Status \\
			\hline
			1 (Isotropie-Lemma)        & \texttt{isotropy}     & PASS \\
			2 (Projektor-Idempotenz)   & \texttt{idempotence}  & PASS \\
			3 (Dualitätsstabilität)    & \texttt{duality}      & PASS \\
			4 (Volumen-Kompatibilität) & \texttt{volume}       & PASS \\
			\hline
			\multicolumn{2}{r}{Gesamt:} & 4/4 PASS \\
		\end{tabular}
	\end{center}
	
	%=============================================================================
	\section*{Synthese und Schlussfolgerung}
	\label{sec:conclusion}
	%=============================================================================
	
	\begin{center}
		\label{box:conclusion}
		\fbox{%
			\fbox{%
				\begin{minipage}{0.92\linewidth}
					\vspace{2mm}
					\begin{center}
						\textbf{\large Synthese der EABC-Renormierung}
					\end{center}
					\vspace{3mm}
					\begin{equation*}
						\textbf{Hauptsatz (24-Isotropy-Restoration, Theorem~\ref{thm:tensor-restore})}\!:
						\quad
						M_{\mathrm{eff}}\bigl(R^*_{\mathrm{EABC}}(K^+)\bigr) = 24I_3.
					\end{equation*}
					\begin{equation*}
						\textbf{Vorbereitung}\!:
						\quad
						M_{\mathrm{geom}}=24I_3,\;
						M_{\mathrm{eff}}(K^+)=24I_3+w_pvv^{\mathsf T},\;
						\Delta(M_{\mathrm{eff}})=w_p.
					\end{equation*}
					\begin{equation*}
						\textbf{Peano-Fortsetzung (Korollar~\ref{cor:shellstack-global})}\!:
						\quad
						\forall i:\;
						M_{\mathrm{eff}}\!\bigl(\mathcal{R}^*(\mathcal{S}_{\mathrm{def}})\bigr)_i = 24I_3.
					\end{equation*}
					\vspace{2mm}
					\hrule
					\vspace{2mm}
					\small
					\textbf{Kernaussage.}
					Eine EABC-Primzahl $p>3$ mit $p\bmod 12\in\{1,5,7,11\}$ induziert auf einer
					Normschale eine quantisierte, messbare Rang-1-Anisotropie
					$w_p\in\{1,5,7,11\}$.
					Der komponentenweise Retraktions-Schnitt $\mathcal{R}^*$ führt das
					Schalensystem algebraisch stabil auf den isotropen Ikosaeder-Fixpunkt zurück.
					Die Satzkette ist in \texttt{eabc-renorm} maschinell verifiziert
					(Lean~4.29, \texttt{lake build} grün, \texttt{sorry}\,$=\,0$;
					Tabelle~\ref{tab:lean-theorems}).
					\textbf{Phase~A} ist der Hauptsatz (Theorem~\ref{thm:tensor-restore}).
					\textbf{Phase~B} liefert kombinatorische Einordnung (eingefroren).
					\textbf{Phase~C} hat einen Teilstand ($n\le 3$, Ikosaeder, Dirac-Anker);
					Box-Dimension und Grenzwert $n\to\infty$ bleiben offen.
					\smallskip
					
					\textit{Intellektuelle Linie:}
					Von Eulers polyedrischer Symmetrie über Riemanns Spektralgeometrie und
					Diracs gestörtem Vakuum zu 't~Hoofts projektiver Renormierungslogik ---
					hier als diskrete, maschinengeprüfte Satzkette auf dem Ikosaeder-Fixpunkt.
					\vspace{2mm}
				\end{minipage}%
			}%
		}
	\end{center}
	
	%=============================================================================
	\appendix
	\section{Explizite Erzeuger-Verifikation}
	\label{app:generators}
	%=============================================================================
	
	\subsection{Ikosaeder-Ecken und geometrische EABC-Zuweisung}
	
	Mit $\varphi=(1+\sqrt5)/2$ seien die zwölf normierten Ikosaeder-Ecken
	\[
	v_i \in
	\Bigl\{
	\frac{(0,\pm1,\pm\varphi)}{\|(0,1,\varphi)\|},
	\frac{(\pm1,\pm\varphi,0)}{\|(1,\varphi,0)\|},
	\frac{(\pm\varphi,0,\pm1)}{\|(\varphi,0,1)\|}
	\Bigr\}.
	\]
	Die Koordinaten-Restklasse $r(v)$ mit Skala $s=6$ teilt die Ecken in vier
	Bahnen zu je drei Elementen mit Restklassen $2,4,8,10\bmod 12$.
	Die Abbildung \emph{nächste EABC-Restklasse} liefert
	\[
	2\mapsto E\;(w=1),\quad
	4\mapsto A\;(w=5),\quad
	8\mapsto B\;(w=7),\quad
	10\mapsto C\;(w=11).
	\]
	
	\subsection{Direkte Berechnung von $M$}
	
	Am geometrischen Fixpunkt $I_{12}^{\mathrm{EABC}}$ gilt
	\[
	M
	\;=\;
	\sum_{i=1}^{12} w_{\sigma(i)}\, v_i v_i^{\mathsf T}
	\;\in\;
	\Q(\sqrt5)^{3\times 3}.
	\]
	Eine direkte Summation über die zwölf äußeren Produkte (symbolisch in
	$\varphi$) ergibt exakt
	\[
	M = 24\,I_3.
	\]
	Damit ist $\Delta(M)=0$ und $\lambda_1=\lambda_2=\lambda_3=24$ unmittelbar
	sichtbar, unabhängig von der Symmetrie-Rechnung im Beweis von Lemma~1.
	
	\subsection{Erzeuger $R_3$ und $R_5$}
	
	Die $2\pi/3$-Drehung um den Flächenschwerpunkt $c=\tfrac{1}{\sqrt3}(1,1,1)$
	hat die Matrix
	\[
	R_3
	\;=\;
	\begin{pmatrix}
		0&1&0\\ 0&0&1\\ 1&0&0
	\end{pmatrix}
	\;\in\; SO(3).
	\]
	Die $2\pi/5$-Drehung um die Ecke $v_0$ (Rodrigues-Formel) liefert
	\[
	R_5
	\;=\;
	\begin{pmatrix}
		\tfrac{-1+\varphi}{2} & \tfrac{1+\varphi}{2} & -\tfrac12 \\[4pt]
		-\tfrac{1+\varphi}{2} & \tfrac12 & \tfrac{-1+\varphi}{2} \\[4pt]
		\tfrac12 & \tfrac{-1+\varphi}{2} & \tfrac{1+\varphi}{2}
	\end{pmatrix}
	\;\in\; SO(3),
	\]
	wobei man $\cos(2\pi/5)=\tfrac{-1+\varphi}{2}$ und
	$\sin(2\pi/5)=\tfrac{\sqrt{10+2\sqrt5}}{4}$ verwendet.
	
	\subsection{Verifikation der $\mathcal{I}$-Invarianz}
	
	Da $M=24I_3$ exakt gilt, folgt für jedes $R\in SO(3)$ sofort
	\[
	R\,M\,R^{\mathsf T}
	\;=\;
	24\,R\,R^{\mathsf T}
	\;=\;
	24\,I_3
	\;=\;
	M.
	\]
	Insbesondere für die Erzeuger $R_3,R_5\in A_5$:
	\[
	R_k\,M\,R_k^{\mathsf T} = M,
	\qquad k=3,5.
	\]
	Damit ist Schritt~(iii) im Beweis des Isotropie-Lemmas vollständig explizit
	nachgerechnet; die Untergruppe $\mathcal{G}\subset\mathcal{I}$ aus dem
	Hauptbeweis enthält $R_3,R_5$ und daher ganz $\mathcal{I}\cong A_5$.
	
	\bibliographystyle{plain}
	\bibliography{references}
	
\end{document}
